\documentclass[11pt,a4paper]{article}
\usepackage[T1]{fontenc}
\usepackage[margin=27mm]{geometry}
\usepackage{amsmath,amssymb,amsthm,mathtools}
\usepackage{booktabs,array,longtable}
\usepackage[expansion=false]{microtype}
\usepackage[hidelinks]{hyperref}
\usepackage{url}
\newcolumntype{P}[1]{>{\raggedright\arraybackslash}p{#1}}
\setlength{\emergencystretch}{3em}
\allowdisplaybreaks[2]
\numberwithin{equation}{section}
\newtheorem{theorem}{Theorem}[section]
\newtheorem{lemma}[theorem]{Lemma}
\newtheorem{corollary}[theorem]{Corollary}
\theoremstyle{definition}
\newtheorem{definition}[theorem]{Definition}
\newtheorem{hypothesis}[theorem]{Hypothesis}
\theoremstyle{remark}
\newtheorem{remark}[theorem]{Remark}
\newcommand{\HTC}{\ensuremath{H_{\mathrm{TC}}}}
\newcommand{\HBB}{\ensuremath{H_{\mathrm{BB}}}}
\newcommand{\poly}{\operatorname{poly}}
\newcommand{\THR}{\mathrm{THR}}
\newcommand{\ETHR}{\mathrm{exact\text{-}THR}}
\newcommand{\MOD}{\mathrm{MOD}}
\newcommand{\ENP}{\mathrm{E}^{\mathrm{NP}}}
\newcommand{\NTIME}{\mathrm{NTIME}}
\newcommand{\cc}{\operatorname{cc}}
\newcommand{\inv}{\operatorname{inv}}
\newcommand{\sgn}{\operatorname{sgn}}
\newcommand{\Enc}{\operatorname{Enc}}
\newcommand{\Dec}{\operatorname{Dec}}
\newcommand{\bits}{\operatorname{bits}}
\newcommand{\E}{\mathbb E}
\newcommand{\Z}{\mathbb Z}
\newcommand{\Q}{\mathbb Q}
\newcommand{\eps}{\varepsilon}
\newcommand{\ind}[1]{\left[#1\right]}
\newcommand{\ancfile}[1]{\texttt{\detokenize{#1}}}
\newcommand{\gfun}[1]{\frac{#1}{\sqrt{\log_2(#1+16)}}}
\hypersetup{pdftitle={Conditional consequences of a Boolean tensor-contraction theorem},pdfauthor={Jizhou Guo}}
\title{Conditional consequences of a Boolean\\ tensor-contraction theorem:\\ threshold circuits and sparse connectivity}
\author{Jizhou Guo\\\texttt{mitsuha2021b@gmail.com}}
\date{29 October 2026}
\begin{document}
\maketitle
\begin{abstract}
Assuming the Boolean tensor-contraction theorem in OpenAI's manuscript
\emph{The Exponential Time Hypothesis is False} (25 October 2026), we collect
consequences for threshold circuits and sparse graph problems.
For an $n$-input circuit with $w$ wires and a polynomial length integer
weight description, put $S=n+w+2$. Its acceptance probability has a
deterministic additive approximation, down to error
$2^{-S^{9/10}/\log_2(S+16)}$, in
$\poly(S)2^{O(S/\sqrt{\log S})}$ time and polynomial space, with unrestricted
depth and fan-in. A per-gate theorem combines exact threshold, equality,
and modular gadgets with hashed gadgets; it gives exact counting when
no gate is hashed. Exact-only assignments allow arbitrary description
length, charged polynomially. Applying established algorithms-to-lower-bounds
constructions yields an $\ENP$ language requiring more than
$cN\sqrt{\log N}$ real-weight threshold wires at every sufficiently large
length, and a corresponding statement for every fixed constant
approximation advantage. For graphs with $m=O(n)$, a sparse representation
of the permutation sign and the matrix-tree identity give deterministic
$2^{O(n/\sqrt{\log n})}$ algorithms for connectivity optimization, witness
recovery, and exact Hamiltonian-cycle counting. Directed and sparse
shortest-path-metric variants are included. Every algorithmic and
complexity-theoretic conclusion is conditional on the stated input
theorem. The accompanying finite computations check the reductions.
\end{abstract}
\tableofcontents
\section{Introduction}

Theorem D.1 of OpenAI's 25 October 2026 manuscript~\cite{OpenAI26}
asserts a deterministic, polynomial-space contraction algorithm for
Boolean tensor networks with linearly many edges and maximum degree
$\lceil s^{9/10}\rceil$. We denote this assertion by \HTC, and state its
full interface in Section~\ref{sec:interface}. This paper studies its
consequences. The degree bound, the allowed coefficient ring, the cost
of evaluating implicit tensors, and exact integer readback are part of
every application.

Throughout, $\log_2$ is the binary logarithm, and
$g(s)=s/\sqrt{\log_2(s+16)}$. Running times count bit operations, including
polynomial factors in the description length. Polynomial space means
polynomially many bits. A circuit wire is an occurrence in a gate's
input list; arbitrary fan-out is permitted. Circuit depth and fan-in
are unrestricted unless a statement explicitly bounds them.

\subsection{Statements of the consequences}

\begin{theorem}[Threshold acceptance probability, conditional on \HTC]
\label{thm:capp}
Assume \HTC. Fix $a\ge1$. Let $C$ be a single-output Boolean threshold
circuit with $n$ declared inputs, $w$ wires, and integer weights and
thresholds in a complete description of length $L\le S^a$, where
$S=n+w+2$. For every
\[
 2^{-S^{9/10}/\log_2(S+16)}\le\eps<1,
\]
there is a deterministic algorithm producing a rational $\widehat p$
with
\[
 \left|\widehat p-2^{-n}|\{x:C(x)=1\}|\right|\le\eps
\]
in $\poly_a(S+L)2^{O_a(g(S))}$ time and $\poly_a(S+L)$ space.
In particular this holds for $\eps\ge S^{-a}$. The algorithm can use the
fixed lower endpoint as its target error, so no arbitrary-precision real
encoding of $\eps$ is required.
\end{theorem}

The proof is in Sections~\ref{sec:protocol}--\ref{sec:unified}.
Theorem~\ref{thm:unified} gives the precise per-gate version. It allows
$\THR$, $\ETHR$, and weighted $\MOD_m$ gates, as well as small gates
with polynomial-time evaluators. Each large gate is assigned either an
exact sliced gadget, an exact supplied weight-factorization gadget, or
a hashed gadget. Under the explicit wire budget~\eqref{eq:budget},
the same time and space bounds hold, with error at most the sum of the
hashed gates' error budgets. When that set is empty, the answer is the
exact number of accepting assignments. In that exact-only case
the description length $L$ is arbitrary, with polynomial cost in $S+L$.
In particular, exact counting is available when the wire-weighted sum
of large gates' parameter bit lengths is $O(S^{19/10})$.

\begin{theorem}[Real-weight wire lower bounds, conditional on \HTC]
\label{thm:lower}
Assume \HTC. There are a constant $c>0$ and a language
$L_{\mathrm{ae}}\in\ENP$ such that, for every sufficiently large $N$,
every real-weight threshold circuit computing
$(L_{\mathrm{ae}})_N$ has more than $cN\sqrt{\log_2N}$ wires.
Consequently there is also such a language with this lower bound at
infinitely many lengths. For every fixed $\gamma\in(0,1/2)$ there are
$c_\gamma>0$ and $L_\gamma\in\ENP$ such that, for every sufficiently
large $N$, every real-weight threshold circuit with at most
$c_\gamma N\sqrt{\log_2N}$ wires agrees with $(L_\gamma)_N$ on at most
a $1/2+\gamma$ fraction of inputs. The circuits have arbitrary depth
and fan-in. The language in the last assertion may depend on $\gamma$.
\end{theorem}

Section~\ref{sec:lower} proves these statements by substituting the
acceptance-probability algorithm into the proofs of known transfer
constructions. We also give the separate infinitely-often construction,
since it makes the two-stage verification cost explicit.

\begin{theorem}[Undirected sparse connectivity, conditional on \HTC]
\label{thm:undirected}
Assume \HTC. Fix $c>0$. On simple undirected graphs with $n$ vertices
and $m\le cn$ edges, Hamiltonian cycle, Hamiltonian path (with specified
or unspecified endpoints), longest simple path, Hamiltonian TSP,
Steiner tree, feedback vertex set, connected vertex cover, and connected
dominating set admit deterministic decision, optimization, and witness
recovery in
\begin{equation}\label{eq:graph-time}
 \poly_c(n+L)2^{O_c(g(n))}\quad\hbox{time and}\quad\poly_c(n+L)
 \quad\hbox{space}.
\end{equation}
The relevant vertex and edge costs may be nonnegative integers of at
most $\lceil n^{9/10}\rceil$ bits. One can also exactly count undirected
Hamiltonian cycles as edge sets, Hamiltonian paths with specified
endpoints, simple paths with specified endpoints and length, and
Steiner tree subgraphs $(U,F)$. For connected covers, connected
dominating sets, and feedback vertex sets, the proof uses positive
weighted counts to decide existence and optimize.
\end{theorem}

\begin{theorem}[Directed sparse connectivity, conditional on \HTC]
\label{thm:directed}
Assume \HTC. Fix $c>0$. On digraphs without loops or parallel arcs and
with $m\le cn$, the bounds~\eqref{eq:graph-time} hold for exact directed
Hamiltonian-cycle counting, exact counting of directed paths with
specified endpoints and length, longest directed simple path,
nonnegative-cost directed Steiner tree, and minimum strongly connected
spanning subgraph. They also hold for a minimum strongly connected
subgraph containing specified terminals. Directed Steiner arborescences
with a fixed root, direction, and terminal set can be counted exactly.
Costs have at most $\lceil n^{9/10}\rceil$ bits. Optimization includes
recovery of an optimal subgraph.
\end{theorem}

\begin{theorem}[Routing on a sparse base graph, conditional on \HTC]
\label{thm:routing}
Assume \HTC. Fix $c>0$. Given a simple undirected base graph with
$m\le cn$, nonnegative integer edge costs of at most
$\lceil n^{9/10}\rceil$ bits, a terminal set $T$, and a required edge
set $R$, a minimum-cost closed walk visiting $T$ and traversing $R$
can be found, or infeasibility detected, within~\eqref{eq:graph-time}.
This includes graphical TSP, Steiner TSP, and undirected Rural Postman.
On a connected base graph, $T=V$ and $R=\varnothing$ give TSP on its
shortest-path metric, with an optimal metric tour recovered for $n\ge3$.
\end{theorem}

\subsection{Comparison and scope}

Table~\ref{tab:comparison} compares exponent scales. In its previous-best
column, $2^{\Theta(n)}$ denotes the single-exponential scale of the
available generic bounds; different problems and graph subclasses have
different constant bases. The cited algorithms specify representative
bounds within that scale. The comparison concerns unrestricted depth in
the circuit rows and a fixed edge-to-vertex ratio in the graph rows.
\begin{table}[ht]
\centering\small
\renewcommand{\arraystretch}{1.18}
\begin{tabular}{@{}P{.29\linewidth}P{.32\linewidth}P{.32\linewidth}@{}}
\toprule
Problem & Previous best (exponent scale) & Under \HTC\\
\midrule
Arbitrary-depth THR CAPP, $w=O(n)$ & $2^{\Theta(n)}$ generic enumeration & $2^{O(g(n))}$, error down to $2^{-\Theta(n^{.9}/\log n)}$\\
Mixed-gate exact counting under~\eqref{eq:budget} & $2^n\poly(L)$ enumeration & $\poly(S+L)2^{O(g(S))}$ for arbitrary $L$ if no gate is hashed\\
Unrestricted-depth real THR wires for explicit $\ENP$ languages & Linear scale; depth-restricted bounds differ & $>cN\sqrt{\log N}$ a.e.; fixed-advantage average case\\
Sparse Hamiltonicity & Single exponential; e.g. randomized $O^*(1.657^n)$~\cite{Bjorklund14} & Deterministic $2^{O(g(n))}$; exact \#HC\\
Sparse paths, Steiner, FVS, connected cover/domination & $2^{\Theta(n)}$ scale; connectivity frameworks~\cite{CutCount11,BCKN15} & Deterministic $2^{O(g(n))}$, optimization and witnesses\\
Sparse directed connectivity and sparse-base routing & $2^{\Theta(n)}$ scale (enumerate selected arcs or edge multiplicities) & Deterministic $2^{O(g(n))}$\\
\bottomrule
\end{tabular}
\caption{Polynomial factors in the input length are suppressed.
Algorithmic comparisons are deterministic unless randomization is marked. The
conditional bounds retain polynomial space. The exact-counting rows use
the object conventions in the theorem statements.}
\label{tab:comparison}
\end{table}

The polynomial description condition in Theorem~\ref{thm:capp} is a bit
length condition. For lower bounds, each real-weight threshold gate has
an equivalent polynomial-bit integer representation, with no change
in wires. This is an existence assertion, sufficient for guessing
circuits; it does not prescribe how to read arbitrary real-number
oracles. For an algorithm on a larger wire budget, the displayed bound
is in $S$, and may exceed exhaustive enumeration over the $n$ inputs.

\subsection{Related work}

OpenAI's manuscript itself derives a circuit lower bound from the
fixed-density case of its Theorem~1.1~\cite[Corollary~1.4]{OpenAI26}:
there are $L\in\ENP$ and $c>0$ such that circuits of fan-in-two Boolean
gates computing $L$ need more than $cn\sqrt{\log n}$ noninput gates for
infinitely many $n$. Its proof combines Williams' witness-circuit
argument with a local reduction to 3-CNF. Theorem~\ref{thm:lower} uses the
stronger input \HTC\ and a different measure. Every fan-in-two Boolean
gate has a threshold implementation with at most five wires, so
Theorem~\ref{thm:lower} implies a gate lower bound of the same order for
fan-in-two circuits, at every sufficiently large length and in the
constant-advantage form. In the other direction, the generic simulation
of Section~\ref{sec:lower} replaces a fan-in-$r$ threshold gate by
$O(r^2\log(r+2))$ fan-in-two gates, so Corollary~1.4 does not yield the
wire bound of Theorem~\ref{thm:lower} by gate simulation.

Williams' algorithmic method~\cite{Williams13} connects savings in
circuit analysis to circuit lower bounds. Chen and Williams~\cite{CW19}
use PCPs of proximity to reduce the composition overhead in such
connections. Bathie and Williams~\cite[Section 3.2--3.3]{BW24} give a
two-stage PCP/ECC/PCPP construction whose circuit queries are small
compositions of proof-bit circuits. The infinitely-often argument here
is a routine instantiation of that proof framework in the wire measure.
Chen, Lyu, and Williams~\cite{CLW20} give the refuter for a hierarchy
with sublinear nondeterminism and an XOR amplification based on linear
sums. Their proof mechanisms supply the almost-everywhere and
constant-advantage arguments below; the polynomial parameters and the
per-term wire counts are recalculated explicitly.

Depth-restricted threshold bounds have different quantitative behavior.
Impagliazzo, Paturi, and Saks~\cite{IPS97} prove that parity requires
$n^{1+c\theta^{-d}}$ wires at depth $d$, for absolute constants $c>0$
and $\theta\le3$. Kane and Williams~\cite{KW16} prove, for a fixed
advantage, depth-two lower bounds of order
$n^{5/2}/\log^{7/2}n$ wires and $n^{3/2}/\log^3n$ gates for an explicit
function; their depth-three result has a majority output gate.
Chen, Tal, and Wang's 2026 manuscript~\cite[Theorems 1.1--1.2]{CTW26}
states $n^{2.5-\eps}$ gate lower bounds in $\ENP$ for depth-two linear
threshold circuits, via a CAPP algorithm for an XOR of two such
circuits. These gate and wire measures and their depth restrictions
are retained in the comparison. Li and Yang~\cite{LiYang22} prove a
$3.1n-o(n)$ gate lower bound over the full binary Boolean basis. A
binary Boolean gate has a constant-size threshold implementation; the
reverse simulation of an arbitrary-weight threshold gate has a larger
bit cost, as quantified in Section~\ref{sec:lower}.

The comparison protocol used here follows Nisan's
integer-comparison construction~\cite{Nisan93}. Viola~\cite{Viola15}
improves the communication for addition and comparison. We use an
explicit half-splitting and prime-fingerprinting implementation to keep
seed length, local evaluation, and tensor ports in the same ledger.
The conversion from a deterministic transcript to local consistency
tensors and the union bound over the earliest erroneous gate are standard constructions.

For graph connectivity, Cut\&Count~\cite{CutCount11} combines consistent
cuts with isolation to filter disconnected candidates modulo two.
Bodlaender, Cygan, Kratsch, and Nederlof~\cite[Section 4]{BCKN15}
use determinants and the matrix-tree theorem for deterministic
connectivity counting parameterized by treewidth. Bj\"orklund's
Hamiltonicity algorithm~\cite{Bjorklund14} uses a different determinant
summation and cancellation. Our graph proofs instantiate the
matrix-tree approach using an explicit representation of the sparse
permutation sign by linearly many Boolean wires. The resulting scale is
$n+m$; no treewidth parameter is substituted for the number of vertices.
The $0/1/2$ multiplicity reduction for graphical TSP also appears in
the Cut\&Count treatment~\cite[Section 4.10]{CutCount11}.

\section{The contraction interface and encoding conventions}
\label{sec:interface}

\begin{hypothesis}[\HTC: OpenAI, Theorem D.1]
Fix positive constants $C_0,A_0,b_0$. For an integer $s\ge2$, let a
bipartite multigraph carry one Boolean variable on each edge. There
are at most $C_0s$ edges and a number of vertices bounded by a fixed
polynomial in $s$, on which the constants below may also depend, and its
maximum degree is at most $\lceil s^{9/10}\rceil$. If the total graph
and tensor description length is $\mathcal L$, the tensor entries lie
in
\[
 \mathcal B_{P_0,\ell}=(\Z/2^{P_0}\Z)[V]/(V^\ell),
 \qquad 1\le P_0,\ell\le C_0s.
\]
Uniformly for every vertex and every Boolean port assignment, an entry
can be computed from its description in time
$(s+\mathcal L)^{b_0}2^{A_0s^{9/10}}$ and space
$(s+\mathcal L)^{b_0}$. Then the contraction
\begin{equation}\label{eq:contraction}
 Z=\sum_{x\in\{0,1\}^{E}}\prod_v T_v(x|_{\delta(v)})
\end{equation}
is computable deterministically in
$(s+\mathcal L)^{O(1)}2^{O(g(s))}$ bit operations and
$(s+\mathcal L)^{O(1)}$ space. Constants depend only on the fixed
bounds just specified. Isolated vertices contribute their scalar
entries; $\ell=1$ is allowed.
\end{hypothesis}

This is Theorem~D.1 (Boolean tensor contraction) of~\cite{OpenAI26}.
All our applications use $\ell=1$. The only
coefficients required are integers modulo a power of two, including
negative residues for determinants.

For comparison, write \HBB\ for Theorem 1.1 of~\cite{OpenAI26}:
for each fixed $k\ge2$, exact counting for $k$-CNF on $n\ge20$
declared variables and $m$ clauses takes polynomial space and time
\[
 (n+m)^{O_k(1)}
 2^{O_k(n\log\log n/\sqrt{\log n})}.
\]
For each fixed $D$, if $m\le Dn$, the exponent improves to
$O_{k,D}(g(n))$. The quantifier fixes $k$ and $D$ before the input
size grows.

The accompanying Lean scope document~\cite{OpenAILean26},
\ancfile{lean/docs/378.md}, describes exact counting for fixed $k$ on
canonical encodings, in time
\[
 (L+1)^{O_k(1)}
 2^{O_k(n(\log\log n)^2/\sqrt{\log n})}
\]
and with polynomially many visited tape cells. Its stated scope is
this fixed-$k$ counting assertion. It does not include \HTC, nor the
sharper arbitrary-density or fixed-density bounds in the October 25
Theorem 1.1. Thus that formalization does not certify the hypothesis
used by the results of this paper.

\subsection{Boolean wires and unique extensions}

A functional gadget has input and output ports and internal Boolean
wires. For each assignment to its input ports, exactly one assignment
to its internal wires and output ports has nonzero product of local
$0/1$ entries. A circuit of such gadgets has a unique extension of each
original input, by induction in a topological order. Pinning the final
output to one then makes its contraction the accepting-input count.

Fan-out is handled by equality trees. For a signal with $d\ge2$
occurrences, join its ports by a tree of binary and ternary equality
vertices; fixing the common signal determines every internal bit.
For $d=1$, attach the free unary tensor $(1,1)$. A declared input
with no occurrence is represented by a wire between two copies of
$(1,1)$, contributing exactly two. In particular, changing a gate to a
constant never deletes the factor contributed by an input that thereby
becomes unused. These rules also apply when the circuit output is an
input or a constant.

A constant gate is evaluated directly and its outgoing uses are fixed
by unary pins. A nonconstant gate has at least one input wire, so the
number of remaining gates is at most $w$. Constants, dead gates, and
zero-arity scalar factors can be removed from the network description
after accounting for their effect and retaining all declared inputs.

Given any finite factor graph, insert a binary equality tensor on
each edge. The old and new vertices form a bipartition, edge count
doubles, maximum degree does not increase, and the assignment on each
old edge has exactly one extension to its two replacements. If the
sum of port occurrences is $O(S)$, equality trees and subdivision
therefore use $O(S)$ edges and vertices.
Distinct occurrences at the same factor are distinct ports. If a
two-occurrence signal joins two ports of that factor, subdivision
creates parallel edges to the equality vertex. These are allowed by
the bipartite multigraph interface of \HTC.

\subsection{Scaling and integer readback}

Suppose an encoding has $C_1S$ edges, $O(S)$ vertices, degree at most
$D_1S^{9/10}$, precision $P_0\le C_2S$, and total description length
$\mathcal L$. Suppose its local entry programs use fixed polynomial
time and space in $S+\mathcal L$. Set
\begin{equation}\label{eq:rescale}
 s=\lceil KS\rceil,\qquad
 K\ge\max\{1,(2D_1)^{10/9}\},
\end{equation}
and fix a sufficiently large interface constant $C_0$. All the bounds
in \HTC\ now hold, including its precise maximum-degree condition.
Since $K$ is fixed, $g(s)=O(g(S))$. The same reasoning applies to
$O(S^{2/3})$ degrees. The resulting time is
$(S+\mathcal L)^{O(1)}2^{O(g(S))}$ and the space is polynomial in
$S+\mathcal L$, allowing arbitrary description length. Bounded small
instances may be enumerated with polynomial cost in their descriptions.

If a contraction is known, as an integer, to satisfy
$0\le Z<2^{P_0}$, its least nonnegative residue modulo $2^{P_0}$
recovers $Z$ exactly. For a circuit with unique extensions, $P_0=n+1$
is sufficient, since $0\le Z\le2^n$. For a signed graph encoding the
bound is proved after cancellation, before choosing $P_0$. Division by
a known multiplicity is performed on this recovered integer.

\part{Threshold circuits}
\section{Comparison protocols as bounded-degree relations}
\label{sec:protocol}

We give the comparison implementation used in the threshold gadgets.
It is the half-splitting comparison underlying Nisan's
protocol~\cite{Nisan93}; explicit parameters are useful for applying
\HTC. Fix the circuit scale $S$, and suppose $t\le S$ parties hold
nonnegative integers $X_i$, each of bit length polynomial in $S$.
The public target $T$ has the same bit bound. We compare
$\sum_iX_i>T$. Targets outside the possible sum range are handled by
public constant answers. Pad the bit width to a power of two $B'$ and
put $J=\log_2 B'=O_a(\log S)$ when the input description is at most
$S^a$ bits.

\subsection{The half-splitting step}

For a current even width $B$, put $d=B/2$ and write
\[
 X_i=2^dH_i+L_i,\quad T=2^dT_H+T_L,
 \quad 0\le L_i,T_L<2^d.
\]
If $\sum_iH_i=T_H-h$ for some $h\in\{0,\ldots,t-1\}$, then the
comparison is equivalent to
\begin{equation}\label{eq:low-comparison}
 \sum_i L_i>T_L+h2^d.
\end{equation}
If there is no such $h$, put $u=\sum_iH_i-T_H$. Then either $u\ge1$
or $u\le-t$. In the former case the original comparison is true, since
$2^du+\sum_iL_i-T_L\ge1$; in the latter case it is false. Both
answers equal the answer to $\sum_iH_i>T_H$. Thus one may recurse on
that high-half comparison. The range for $h$ includes $t-1$.

Let $\delta=2^{-h_0}$, with integer $h_0\ge1$. Choose the least power
of two $K$ such that
\begin{equation}\label{eq:prime-family}
 K\ge 4JS(B'+\lceil\log_2S\rceil+3)\,2^{h_0}.
\end{equation}
Use the $K$ smallest primes greater than $S$. At each level a fresh
$\log_2K$-bit seed chooses one such prime $p$. The parties broadcast
$H_i\bmod p$, and public arithmetic computes the least nonnegative
residue
\[
 h'=T_H-\sum_i(H_i\bmod p)\pmod p.
\]
If $h'<t$, follow~\eqref{eq:low-comparison} with $h=h'$; otherwise
follow the high-half comparison. Since $p>t$, a genuine $h$ is always
returned correctly. When no genuine $h$ exists, an incorrect choice
requires $p$ to divide one of the nonzero integers
\[
 \sum_iH_i-T_H+j,\qquad 0\le j<t.
\]
After the range checks, each has bit length at most
$B'+\lceil\log_2S\rceil+3$. An integer of that bit length has at most
that many distinct prime factors. Consequently the conditional error
at a level is at most $\delta/(4J)$, even after fixing all earlier
messages and seeds. A union bound on the earliest erroneous level gives
error at most $\delta$. At width one the parties send their remaining
bits and compare exactly.

The prime bound of Rosser and Schoenfeld~\cite{Rosser62} implies that
the largest prime used has
\begin{equation}\label{eq:prime-bits}
 \pi=O_a(h_0+\log S)
\end{equation}
bits. Each seed is only a prime index. For a particular index,
enumerate integers and use trial division until reaching that prime.
The polynomial upper bound on the indexed prime implies
$2^{O_a(h_0+\log S)}$ time and polynomial bit space. This construction
stores the selected primes, rather than the entire family. The total
seed length is
\begin{equation}\label{eq:seedlength}
 R=J\log_2K=O_a(h_0\log S+\log^2S).
\end{equation}
Fixed-width message frames, range checks, and forced zero padding
after an early stop make every transcript canonical.

\subsection{Two tensor realizations}

Fix the seed. If a party reads at most $b$ original input bits and
the full transcript has $C$ bits, make one factor for each party,
reading its private input and a copy of the transcript. The factor
checks only that party's messages, each against the earlier prefix.
A referee factor checks public control, formatting, padding, and the
output. Sharing transcript bits by equality trees gives
\begin{equation}\label{eq:full-transcript}
 M=O(tb+tC+C+1),\qquad
 \Delta\le\max\{b+C,C+1,3\}.
\end{equation}
For fixed private inputs, the initial message is determined. Inductively
each next message, stop flag, and final output is determined; the
checks force exactly this execution. Hence the gadget is a functional
relation even for a seed at which the comparison is wrong.

For higher precision we use the addition-specific compact
realization. Encode each public branch (a value of $h$ or ``no such
value'') in $\lambda=\lceil\log_2(t+1)\rceil$ bits. The factor for
party $i$ reads its original input, the public branch history, its own
$J$ residues, and its final remaining bit. It replays its local
state and checks these messages. It reads no other party's residues.
At each level a public factor uses earlier history to produce the
target residue and the stop flag. A fan-in-two Boolean circuit adds
the $t$ residues modulo the known prime, subtracts from the target
modulo that prime, and compares the result with $t$. Addition of two
$\pi$-bit residues, comparison with $p$, subtraction, and multiplexing
have $O(\pi)$ gates and wires. Therefore this public arithmetic costs
$O(t\pi)$ wires per level. A final factor checks the output from the
$t$ last bits and the public history.

Counting original input ports, residues, arithmetic wires, and history
copies gives
\begin{align}
 M&=O\bigl(tb+tJ\pi+tJ^2\lambda+J\pi+1\bigr),
       \label{eq:compact-edges}\\
 \Delta&=O\bigl(\max\{b+J(\pi+\lambda),\,
 \pi+J\lambda,\,t+J\lambda,\,3\}\bigr).
       \label{eq:compact-degree}
\end{align}
The history-copy bound is a convenient overestimate of the actual
port count. Given the inputs, earlier branches determine current
messages; those messages determine the next branch. This induction
also proves uniqueness for the compact realization. All local entry
programs use polynomial-time integer arithmetic on the original
weights and the selected primes.

\subsection{Shared seeds and the inverse-polynomial CAPP bound}

For a circuit, suppose gate $v$ is replaced by $f_{v,\rho}$ satisfying
\[
 \Pr_\rho[f_{v,\rho}(y)\ne f_v(y)]\le\eta_v
 \quad\hbox{for each fixed local input }y.
\]
Every gate may read a prefix of the same uniform seed. For an original
input $x$, evaluate the original circuit and let $y_v(x)$ be its fixed
local input at gate $v$. If the replaced circuit disagrees, its earliest
erroneous gate has input exactly $y_v(x)$. Therefore
\begin{equation}\label{eq:gate-error}
 \Pr_\rho[C_\rho(x)\ne C(x)]
 \le\sum_v\Pr_\rho[f_{v,\rho}(y_v(x))\ne f_v(y_v(x))]
 \le\sum_v\eta_v.
\end{equation}
No independence between gates is needed.

For clarity, the inverse-polynomial case already follows from the
full-transcript realization. Put $b=\lceil S^{2/3}\rceil$ and leave
smaller gates exact. For a large threshold gate, replace negative
weights by literals: with $u_j=|w_j|$ and
$\xi_j=x_j$ for $w_j\ge0$, $\xi_j=1-x_j$ otherwise,
\[
 \sum_jw_jx_j\ge\theta
 \quad\Longleftrightarrow\quad
 \sum_j u_j\xi_j>\theta+\sum_{w_j<0}|w_j|-1.
\]
Party $i$ computes the nonnegative sum in one block of at most $b$
inputs. Give each large gate error $\eta=\eps/(2S)$, using the
protocol with a dyadic error at most this value. If $\eps\ge S^{-a}$,
then $R=O_a(\log^2S)$ and a gate with fan-in $r_v$ has
$t_v\le2r_v/b$ parties and transcript
$C_v=O_a(t_v\log^2S)$. The total additional edges in
\eqref{eq:full-transcript} satisfy
\[
 \sum_{r_v>b}O(t_vC_v)
 \le O_a\left(\frac{\log^2S}{b^2}\sum_vr_v^2\right)
 \le O_a(S^{2/3}\log^2S)=o(S).
\]
The degree is at most
$S^{2/3}+O_a(S^{1/3}\log^2S)+O(1)=o(S^{9/10})$.
For each fixed seed, the circuit has a unique-extension network and
hence an exact count $Z_\rho$ from \HTC\ with $P_0=n+1$.
Enumerating all seeds gives the rational
\begin{equation}\label{eq:capp-output}
 \widehat p=\frac{\sum_\rho Z_\rho}{2^{n+R}}.
\end{equation}
There are at most $S$ gates, so~\eqref{eq:gate-error} bounds the
error by $\eps$. The seed enumeration is absorbed by
$2^{O_a(g(S))}$, and the accumulator has $n+R+1$ bits. This proves
the inverse-polynomial assertion of Theorem~\ref{thm:capp} under
\HTC. The compact realization supplies its stated endpoint in the
next section.

\section{Per-gate exact and approximate counting}
\label{sec:unified}

Let $C$ be a single-output Boolean DAG. Its gates have ordered input
lists, possibly with repetitions. Put $S=n+w+2$ and
$b=\lceil S^{9/10}\rceil$. A gate of fan-in $r_v>b$ is large; all
other gates are small. The allowed types are
\[
\begin{aligned}
\THR &: \ind{\sum_jw_jy_j\ge\theta},\\
\ETHR &: \ind{\sum_jw_jy_j=\theta},\\
\MOD_m &: \ind{\sum_jw_jy_j\equiv\theta\pmod m},
\end{aligned}
\]
where all parameters are integers and $m\ge1$. Let $L$ be the length
of the complete input description, including the mode assignment and
the certificates below, and fix $a\ge1$. A small gate can instead
be an arbitrary Boolean function supplied by an evaluator. When any
gate is hashed, require $L\le S^a$ and evaluator time at most $S^a$.
For an exact-only assignment, $L$ may be arbitrary and evaluator time
at most $(S+L)^a$ is sufficient.

For a large threshold or exact-threshold gate, let $B_v$ be the maximum
bit length of $|w_j|$, at least one. For a modular gate let
$B_v=\max\{1,\bits(m-1)\}$. Each large gate receives one of the
following modes. In the factorization mode, preprocessing verifies the
dimensions and entry bounds of the certificate and checks each row of
$w=C_vg_v$ by integer arithmetic. A malformed certificate is rejected;
the verification costs polynomial time in $S+L$.
\begin{center}\small
\renewcommand{\arraystretch}{1.15}
\begin{tabular}{@{}P{.16\linewidth}P{.22\linewidth}P{.37\linewidth}P{.15\linewidth}@{}}
\toprule
Mode & Gate types & Cost $\sigma_v$ & Error\\
\midrule
Exact slices & All three & $B_v$ & $0$\\
Exact factorization & All three & $\kappa_v=d_v(\beta_v+\lceil\log_2(r_v+1)\rceil+2)$ & $0$\\
Hashed comparison & $\THR$ & $h_v\lceil\log_2S\rceil$, $h_v\ge1$ & $2^{-h_v}$\\
Hashed equality & $\ETHR$, $\MOD_m$ & $h_v$, $h_v\ge1$ & $2^{-h_v}$\\
\bottomrule
\end{tabular}
\end{center}
For factorization, $C_v$ is an $r_v\times d_v$ integer matrix with
entries of absolute value less than $2^{\beta_v}$, $\beta_v\ge1$,
and $g_v\in\Z^{d_v}$. Its coordinates can be large; their full bit
length is charged to $L$.

\begin{theorem}[Per-gate counting, conditional on \HTC]
\label{thm:unified}
Assume \HTC. Fix $a,c\ge1$. Suppose the assigned modes satisfy
\begin{equation}\label{eq:budget}
 \sum_{v\text{ large}}r_v\left\lceil\frac{\sigma_v}{b}\right\rceil
 \le cS,
 \qquad \sigma_v\le cb
 \quad\text{for every factorized or hashed large gate}.
\end{equation}
Let $\mathcal H$ be the set of hashed gates. If
$\mathcal H\ne\varnothing$, under $L\le S^a$, in
$\poly_{a,c}(S)2^{O_{a,c}(g(S))}$ time and polynomial space a
deterministic algorithm outputs a rational $\widehat p$ with
\[
 \left|\widehat p-\Pr_x[C(x)=1]\right|
 \le\sum_{v\in\mathcal H}2^{-h_v}.
\]
If $\mathcal H=\varnothing$, the restriction $L\le S^a$ is unnecessary:
the algorithm outputs the exact accepting-input count in
$\poly_{a,c}(S+L)2^{O_{a,c}(g(S))}$ time and
$\poly_{a,c}(S+L)$ space, under the exact-only evaluator condition above.
\end{theorem}

\subsection{Exact sliced gadgets}

A small gate has the single factor $[z=f(y)]$ of degree $r_v+1$.
Now fix a large gate with fan-in $r$, width $B$, and let
$K=\lceil B/b\rceil$. For $\THR$ or $\ETHR$, replace signed weights
by $u_j=|w_j|$ and the corresponding literal $\xi_j$ as above. Write
$T=\theta+\sum_{w_j<0}|w_j|$ and
$\Sigma=\sum_ju_j\xi_j$. We have
$0\le\Sigma<r2^B$. For a threshold gate clamp $T<0$ to zero; for an
exact-threshold gate replace $T<0$ by $r2^B$. In either case clamp
$T>r2^B$ to $r2^B$. These operations preserve the Boolean function.
For a modular gate, replace $w_j$ by $u_j=w_j\bmod m$ and $\theta$
by $\theta'=\theta\bmod m$, with $\xi_j=y_j$. Then $u_j<2^B$ and
$0\le\Sigma<rm$.

Decompose each $u_j$ into base-$2^b$ digits $u_j^{(k)}$, $0\le k<K$.
For each digit position compute
\[
 Y^{(k)}=\sum_ju_j^{(k)}\xi_j,\qquad
 \Sigma=\sum_{k=0}^{K-1}2^{kb}Y^{(k)}.
\]
Partition the $r$ inputs into $t=\lceil r/b\rceil\le2r/b$ blocks.
A chain for each $k$ reads one block at a time and checks the exact
partial-sum update, passing a state of
$\mu=b+\lceil\log_2(r+1)\rceil$ bits. Each input is read $K$ times.
The chain has $O(r)$ wires per slice, degree $b+2\mu$, and unique
states on every original input.

The final arithmetic relates the slice sums by short carries. For
nonnegative constants $A,T'$ occurring below put
\[
 K^*=\left\lceil\frac{B+\lceil\log_2(r+2)\rceil}{b}\right\rceil.
\]
They are less than $2^{K^*b}$. Let $A_k,T'_k$ be their base-$2^b$
digits and $Y^{(k)}=0$ for $k\ge K$. Starting with
$c_{-1}=0,e_{-1}=1$, enforce
\begin{align}
 z_k&=Y^{(k)}+A_k-T'_k+c_{k-1},\nonumber\\
 c_k&=\left\lfloor z_k/2^b\right\rfloor,
 &e_k&=e_{k-1}\wedge[z_k\equiv0\pmod{2^b}].
 \label{eq:carry}
\end{align}
Induction shows
\[
 \Sigma+A-T'=c_{K^*-1}2^{K^*b}+R_0,
 \qquad 0\le R_0<2^{K^*b}.
\]
Thus the difference is nonnegative exactly when the last carry is
nonnegative, and it is zero exactly when that carry is zero and the
last zero flag is one. Every carry lies in $[-1,r+1]$: if the previous
one does, then $-2^b\le z_k\le(r+1)2^b$. An offset encoding therefore
uses $\lceil\log_2(r+3)\rceil$ bits. With $A=0,T'=T$ the last carry
and flag give the required $\THR$ or $\ETHR$ output.

For $\MOD_m$, put $A=m-\theta'$ and introduce a quotient code
$q\in\{0,\ldots,r\}$. Compute two instances of~\eqref{eq:carry},
with $T'=qm$ and $T'=(q+1)m$. The terminal factor requires
\[
 0\le\Sigma+(m-\theta')-qm<m
\]
and outputs one exactly when its left side is zero. There is precisely
one legal $q$, namely $\lfloor(\Sigma+m-\theta')/m\rfloor$; invalid
binary codes are rejected. Every carry factor can read the short
$q$ code and compute its needed digits locally. The same quotient
gadget includes $m=1$ without a separate case.

The slice chains, carries, quotient copies, and output use $O(rK)$
wires, degree $O(b+\log S)$, and polynomial-time local entry programs.
The bound $K^*\le K+1$ holds for sufficiently large $S$; more generally
$K^*=K+O(1)$ suffices. The $O(K\log S)$ carry and quotient wires
are absorbed by $O(Kr)$ because $r>b$. All internal states and the
output are unique once the input is fixed. This establishes the exact
sliced mode.

\subsection{Exact factorization gadgets}

For $w=Cg$, write $\kappa=\kappa_v$ for this gate's state width.
The information to transmit across blocks is a vector,
rather than the weighted sum itself. After block $i$ store
\[
 Y_i=\sum_{j\text{ in blocks }1,\ldots,i}y_jC_{j,\cdot}\in\Z^d.
\]
Each coordinate has absolute value less than $r2^\beta$. A fixed
signed code of $\beta+\lceil\log_2(r+1)\rceil+2$ bits suffices, so
the complete vector uses $\kappa$ bits. Block factors enforce the
updates. The last factor checks the desired comparison, equality,
or congruence of $g\cdot Y_t$ with $\theta$. The supplied identity
$w=Cg$ proves correctness. Updates and output are unique. Before
equality expansion and subdivision, the logical wire count satisfies
\begin{align*}
 r+t\kappa+1
 &\le r+(r/b+1)\kappa+1\\
 &\le r+cr+cb+1\le(2c+2)r,
\end{align*}
using $t=\lceil r/b\rceil\le r/b+1$, $\kappa\le cb$, and $r>b$.
Thus there are $O_c(r)$ wires for every fixed $c\ge1$, and degree at most
$b+2\kappa+O(1)=O_c(b)$ when $\kappa\le cb$.

For example, if a gate's weights have at most $d$ distinct values,
one may take $C$ to be their indicator matrix and $g$ the list of
values. This gives $\beta=1$ and $\kappa=O(d\log S)$ independently
of the bit lengths of those values. The factorization is part of the
input; no search for an optimal one is used.

\subsection{Hashed equality and congruence}

For an exact-threshold gate set
$\Sigma=\sum_jw_jy_j$ and $\mathcal T=\{\theta\}$. For a modular
gate use reduced weights and the target set
\[
 \mathcal T=\{\theta'+qm:0\le q<r\}.
\]
Because $\Sigma\in[0,rm)$, membership in this set is exactly the
original modular predicate. Given a prime $p>S$, replace either
predicate by
\[
 f_p(y)=\ind{\Sigma\bmod p\in\mathcal T\bmod p}.
\]
Every true input remains true. For a false input, a false positive
requires a prime divisor of some nonzero $\Sigma-\tau$,
$\tau\in\mathcal T$. Each such integer has bit length at most
$S^{\lceil a\rceil}+2\lceil\log_2S\rceil+3$. Put
\[
 N_v=r(S^{\lceil a\rceil}+2\lceil\log_2S\rceil+3)
\]
and let $K_v$ be the least power of two at least $N_v2^{h_v}$.
Selecting among the $K_v$ smallest primes greater than $S$ gives error
at most $2^{-h_v}$. Their common bit bound is
$\pi_v=h_v+O_a(\log S)$.

A block chain passes the partial sum modulo $p$ in $\pi_v$ bits.
Its last factor tests membership in the at-most-$r$ reduced targets.
Reject residue codes outside $[0,p)$. For each fixed prime this is a
unique-extension circuit relation, with $O_{a,c}(r)$ wires and degree
$O_{a,c}(b)$ if $h_v\le cb$. In particular the end-state contribution
is bounded by $O_{a,c}(r)$; it need not be at most $r$ when $c>1$.

\subsection{Hashed comparisons}

For a large $\THR$ gate, use the compact comparison realization
\eqref{eq:compact-edges}--\eqref{eq:compact-degree}, grouping $b$
input occurrences per party and choosing error $2^{-h_v}$. Its
parameters satisfy
\[
 J=O_a(\log S),\quad \pi_v=O_a(h_v+\log S),\quad
 \lambda=O(\log S),\quad t\le2r/b.
\]
If $h_v\lceil\log_2S\rceil\le cb$, then
$J\pi_v=O_{a,c}(b)$ and $J^2\lambda=o(b)$. The edge count is
$O_{a,c}(r)$ and the degree is $O_{a,c}(b)$. For the terminal
$J\pi_v$ term, use $r>b$. The seed length is
$O_a(h_v\log S+\log^2S)=O_{a,c}(b)$.
The comparison's pointwise error and the functional property are the
ones proved in Section~\ref{sec:protocol}.

\subsection{Assembling the network}

\begin{proof}[Proof of Theorem~\ref{thm:unified}]
For every hashed gate use the required prefix of a common seed of
length $R$, where $R$ is the maximum of the gate seed lengths and is
zero if no gate is hashed. Both kinds of hashing above give
$R=O_{a,c}(S^{9/10})$. A uniform seed has uniform prefixes, so each
gate retains its own error guarantee. Generate the selected primes
by index, as described above. Per seed the extra time is
$2^{O_{a,c}(S^{9/10})}$ and the space is polynomial.

Replace each gate by its gadget. Connect every occurrence of each
original signal by the equality construction of
Section~\ref{sec:interface}, including a factor two for each unused
input, and pin the final output. Internal signals read by several
factors, including the MOD quotient code and the comparison protocol's
branch history, also use equality trees; these copies are already
included in the gadget port counts. For a fixed seed the resulting
network counts precisely the accepting assignments of $C_\rho$.
Its edges, including equality and bipartite subdivision, are at most
\[
 O_{a,c}\left(n+w+
 \sum_{v\text{ large}}r_v\lceil\sigma_v/b\rceil\right)=O_{a,c}(S).
\]
Its degree is at most $D_{a,c}S^{9/10}$ for a fixed $D_{a,c}$ and its
vertex count is $O_{a,c}(S)$. Its description length and entry costs
are polynomial in $S+L$, hence polynomial in $S$ in the hashed case.
Use~\eqref{eq:rescale} with $s=\lceil K_{a,c}S\rceil$,
$P_0=n+1$, and $\ell=1$. All the hypotheses of \HTC\ hold.

Recover the count $Z_\rho\in[0,2^n]$ exactly. If $R=0$, this is the
answer. Otherwise form~\eqref{eq:capp-output}. Applying
\eqref{eq:gate-error} and averaging over $x$ gives the stated error.
The number of seeds and the prime-generation cost together add
$2^{O_{a,c}(S^{9/10})}$, absorbed by $2^{O_{a,c}(g(S))}$.
All calls are serial, and the accumulator uses $n+R+1$ bits.

In the exact-only case there are no primes to generate and $R=0$.
All arithmetic on weights, thresholds, moduli, and supplied certificates
uses polynomial time and space in $S+L$. The edge, degree, vertex,
and precision bounds above depend on the mode budget, not on an upper
bound for $L$. Retain the original input in the tensor description,
so its length $\mathcal L$ satisfies
$L\le\mathcal L\le\poly_{a,c}(S+L)$. Thus the local evaluation
bound is also polynomial in $S+\mathcal L$, and the same \HTC\ call
gives the stated arbitrary-description bound.
\end{proof}

\begin{proof}[Proof of Theorem~\ref{thm:capp}, including high precision]
Use Theorem~\ref{thm:unified} with the fixed description exponent
$\lceil a\rceil+2$ to include generated mode parameters; all resulting
constants still depend only on $a$.
Assign hashed-comparison mode to every large gate and choose
\[
 h_v=\left\lceil\log_2(2S/\eps)\right\rceil.
\]
Then $\sum_v2^{-h_v}\le\eps/2$, and at the lower endpoint
$h_v\lceil\log_2S\rceil\le2b$ for all sufficiently large $S$.
Each summand in~\eqref{eq:budget} is at most $2r_v$; taking a fixed
$c\ge2$ proves the required budget. For a parameter-free endpoint
implementation one can use the integer target
\[
 h_v=\left\lceil\frac{\lceil S^{9/10}\rceil}
 {\lfloor\log_2(S+16)\rfloor}\right\rceil
 +\lceil\log_2(2S)\rceil+2,
\]
which gives at least the requested accuracy and still obeys the same
constant budget. Finite small $S$ are enumerated. Apply
Theorem~\ref{thm:unified}.
\end{proof}

\begin{corollary}[Exact specializations, conditional on \HTC]
\label{cor:exact}
Assume \HTC\ and the exact-only evaluator condition above; let $L$ be
the arbitrary complete description length.
If all large gates are assigned exact slices and
\[
 \sum_{v\text{ large}}r_vB_v\le c_1S^{19/10}
\]
for a fixed $c_1$, exact counting takes
$\poly_{a,c_1}(S+L)2^{O_{a,c_1}(g(S))}$ time and
$\poly_{a,c_1}(S+L)$ space.
The same holds when every large gate has a supplied factorization
with $\kappa_v\le c_1b$. Small gates may have arbitrary parameters
with their full lengths charged to $L$. In particular, this recovers exact
counting for bounded-fan-in circuits at the allowed degree scale and
for large gates whose weight bit lengths are $O(S^{9/10})$; the
modular sliced mode also allows modulus bit lengths $O(S^{9/10})$.
\end{corollary}
\begin{proof}
For slicing use
$\lceil B_v/b\rceil\le1+B_v/b$ and $\sum_vr_v\le w<S$.
For factorization use
$\lceil\kappa_v/b\rceil\le\lceil c_1\rceil$.
Theorem~\ref{thm:unified} has no hashed gate in either case.
\end{proof}

\subsection{Relations between the specializations}

For integer weights, an equality gate has the identity
\begin{equation}\label{eq:equality-thresholds}
 \ind{\sum_jw_jy_j=\theta}
 =\ind{\sum_jw_jy_j\ge\theta}
 \wedge\neg\ind{\sum_jw_jy_j\ge\theta+1}.
\end{equation}
This uses the same weight vector in two threshold gates and only a
constant-factor increase in wires. Ordinary unweighted MOD gates of
fixed modulus likewise have $O(r)$-wire Boolean, and hence threshold,
simulations: pass the running residue in a constant number of bits.
They therefore fit the threshold-circuit lower-bound framework of
Section~\ref{sec:lower}. For a fixed modulus, arbitrary weights can
also be reduced modulo that modulus before this simulation. The
modular cases requiring the bit-length or certificate accounting here
are predicates with growing modulus and potentially large reduced
weights.

The sliced budget is an aggregate condition even for pure THR
circuits, whereas the factorization condition bounds each supplied
state width $\kappa_v$. For example, a large gate with
$r_v=\Theta(S^{.95})$ and $B_v=\Theta(S^{.94})$ contributes
$\Theta(S^{1.89})$ to the sliced sufficient condition, while its
one-dimensional certificate $C=w$, $g=1$ has width
$\kappa_v=B_v+O(\log S)$, above the factorization cap.
A shorter certificate is a separate property of its weights. Thus
the two sufficient conditions are not identified, including on pure
THR circuits. With $L\le S^a$ they are specializations of the
polynomial-description statement; their arbitrary-$L$ versions use
its exact-only clause.

\section{From acceptance probability to wire lower bounds}
\label{sec:lower}

Write $\THR[s]$ for real-weight, arbitrary-depth threshold circuits
with at most $s$ wires. Fixing or negating an input changes weights
and thresholds without increasing wires; output negation replaces
$[w\cdot y\ge\theta]$ by $[-w\cdot y\ge1-\theta]$ once integer weights
are chosen. Two-input AND, OR, and XOR have constant-wire implementations,
so an $r$-bit parity has $O(r)$ wires.
After a projection fixes every input of a gate, evaluate that gate and
replace it by constant unary tensors, preserving all remaining declared
inputs and adding no wires.

\subsection{Finite descriptions and the transfer ingredients}

A threshold function on $r$ bits has integer weights and threshold of
$O(r\log(r+2))$ bits. Here is the existence proof. Let $D=r+1$ and
for each $x$ put $a_x=(2f(x)-1)(x,-1)$. Move the original threshold
slightly downward to make all margins strict, and scale to obtain
$a_xz\ge1$ for all $x$. The rows have rank $D$, since the rows at
$0,e_1,\ldots,e_r$ do. A feasible point with maximal rank among its
tight rows has $D$ independent tight rows: otherwise move in a
nonzero direction orthogonal to them, choosing its sign to decrease
another inequality, until a new row becomes tight. This
increases the rank. For the resulting invertible matrix $B$, Cramer's
rule applied to $Bz=\mathbf1$ and Hadamard's inequality bound every
numerator and denominator by $D^{D/2}$. Multiplication by
$|\det B|$ gives the asserted integer representation. Constant
functions are handled separately.

Thus a $w$-wire threshold circuit has a normalized description of
$O((n+w)^2\log(n+w+2))$ bits. This remains true after projections.
A normalized fan-in-$r$ gate can also be simulated with
$O(r^2\log(r+2))$ fan-in-two Boolean gates: form the weighted sum
with binary addition and compare. A circuit with
$w=O(t\sqrt{\log t})$ and $n=O(t)$ consequently has ordinary
Boolean size at most $t^3$ for all sufficiently large $t$.

We use the following published ingredients in their specified forms.
They are quoted tools inside the proofs, and do not add an unproved
algorithmic assumption to \HTC.
All fixed probability thresholds used by a machine are chosen rational,
weakening completeness and soundness bounds within their positive gaps
when necessary. Thus the acceptance comparisons have finite encodings.
\begin{enumerate}
\item The unary nondeterministic time hierarchy supplies
$H\in\NTIME(2^t)\setminus\NTIME(2^{t/2})$
\cite[Theorem 2.1 (full version)]{Williams13}.
\item A succinct PCP for such computations has perfect completeness,
constant soundness $\beta<1$, proof and randomness indices padded to
$q=t+O(\log t)$, and polynomial-size query and verification circuits.
For the refuter computation of time $T(n)$ we use the stronger
soundness $n^{-10}$ and $q=\log_2T(n)+O(\log\log T(n))$ from
\cite[Lemma 3.10]{CLW20}. Its queries are projections. The projection
query lists and the verifier's fixed-width CNF give fixed polynomial
bounds in $q$ on the composed verifier size; their construction time
may additionally be polynomial in the original input length.
\item A constant-rate linear error-correcting code $\Enc$ with
polynomial-time decoder $\Dec$ has a fixed decoding radius. The
Circuit-Eval PCPP in \cite[Lemma 24]{CW19}, used in
\cite[Lemma 3.6]{BW24}, has perfect completeness, constant soundness
$\sigma<1$ at a fixed proximity, polynomial proof length and proof
construction, and constraints that are signed ORs of three queried
bits.
\item The refuter in \cite[Theorem 4.6]{CLW20} applies to any
time-constructible $n\le T(n)\le2^{\poly(n)}$. It supplies a language
$A^T_{\mathrm{FS}}\in\NTIME[T(n)]$ and, for each fixed
nondeterministic machine $M$ running in $o(T(n))$ time with at most
$n/10$ guessed bits, finds for every sufficiently large $n$ an input
$z\in\{0,1\}^n$ with $M(z)\ne A^T_{\mathrm{FS}}(z)$, in
$\poly(T(n))$ time with a SAT oracle.
\end{enumerate}

Fix once and for all the algorithm supplied by \HTC; it is part of
the finite description of each cheating machine below. The languages
constructed from their refuters are then fixed languages.
For each fixed description exponent $a$, fix an exponent constant
$A_5(a)$ such that the inverse-polynomial CAPP algorithm takes
\[
 \poly_a(S)2^{A_5(a)g(S)}
\]
time on normalized descriptions of length at most $S^a$, with error
at least $S^{-a}$. This constant is fixed before choosing wire-bound
constants.

All wire caps below can be made computable integer functions. For a
fixed positive rational $d$, use
\begin{equation}\label{eq:cap}
 B_d(r)=\left\lfloor 2d\,\operatorname{isqrt}
 \bigl(r^2\max\{1,\lfloor\log_2r\rfloor\}\bigr)\right\rfloor
 =(2d+o(1))r\sqrt{\log_2r}.
\end{equation}
The floor is computed using rational integer arithmetic. Every
candidate circuit is padded to a predetermined polynomial description
length, and invalid topologies or out-of-range codes are rejected.

\subsection{The infinitely-often construction}
\label{subsec:io}

Assume \HTC. Apply the succinct PCP to the unary hierarchy language
$H$. Fix a common nondecreasing index bound
$q(t)=t+O(\log t)$, and pad proof and randomness indices to this bound.
Extra random bits are ignored and unused proof bits are zero.
The lexicographically least proof accepted on all random strings,
if it exists, defines a proof-bit function $W_1(t,x)$; otherwise set
that function to zero. The full verifier has size $2^{O(t)}$ when all
random strings are written out. Prefix SAT queries recover its
lexicographically least proof in $2^{O(t)}$ time with an NP oracle.

Fix a small positive rational $d$. Define $W_2(t,x,j)$ as follows.
Find the lexicographically least normalized integer threshold circuit
$F_t$ of at most $B_d(q)$ wires whose truth table is a fully accepted
PCP proof. This search is in $2^{O(t)}$ time with SAT queries: for
each candidate description, write out the verification at every one
of the $2^q$ random strings, and regard their conjunction as an NP
condition on the description bits. Let $G_t(x)$ be rejection by the
original PCP using this circuit, and put
\[
 D_t(y)=1-G_t(\Dec(y)).
\]
The verifier $G_t$ is an ordinary Boolean circuit. If an
exact-threshold or MOD gate is used in such a verifier with
polynomial-bit parameters, expand it into binary arithmetic and
comparison circuits of polynomial size before applying the PCPP.
For every $x$ with $D_t(\Enc(x))=1$, take the canonical PCPP proof
from its polynomial-time constructor. Its $j$th bit is $W_2(t,x,j)$.
If no $F_t$ exists, or for an invalid index, output zero. Proof lengths
are padded to one fixed $t^u$ bound covering all legal candidates, so
$j$ has $\ell(t)=O(\log t)$ bits.

Both functions belong to $\ENP$ when inputs include a self-delimiting
$t$ header. Before doing any exponential work, check that $t$ is at
most the full input length and that the payload length is correct.
Form a single language $L_\star$ by a bit selecting $W_1$ or $W_2$.
Its legal lengths are $t+O(\log t)$, with possibly different padding
for the two branches. If $L_\star$ had circuits of at most
$dN\sqrt{\log_2N}$ wires at every sufficiently large $N$, projection
would give circuits for both bit functions within $B_d(q)$.

A fast nondeterministic simulator for $H(1^t)$ guesses $F$ and a
universal proof-bit circuit $T(x,j)$, each of at most $B_d(q)$ wires.
Generate the PCPP constraints $H_1,\ldots,H_m$ of the associated
$D$. Replace every code query by its parity on $x$ and every proof
query by the corresponding fixed-$j$ projection of the \emph{same}
$T$. The resulting threshold circuit $J_i(x)$ satisfies
\begin{equation}\label{eq:three-copy}
 w(J_i)\le3B_d(q)+O(q)
       \le6d\,q\sqrt{\log_2q}+O(q).
\end{equation}
For example, using six wires per two-input XOR gives the explicit
bound $3B_d(q)+18q+3$. The polynomial-size original verifier is used
to generate the PCPP; it is not copied into $J_i$.

If $H(1^t)=1$, the assumed small circuit for $W_1$ makes a legal
$F_t$ exist; the small circuit for $W_2$ supplies $T$, and every
$J_i$ is identically one. If $H(1^t)=0$, every $F$ makes $G(x)=1$
on at least $1-\beta$ of the inputs. For such $x$, every word within
the code's decoding radius of $\Enc(x)$ decodes to $x$, so is rejected
by $D$. The PCPP soundness then makes at least $1-\sigma$ of the
constraints fail for \emph{any} proof at that $x$. Averaging gives
some $i$ with
\[
 \Pr_x[J_i(x)=0]\ge\tau=(1-\beta)(1-\sigma)>0.
\]
CAPP error $1/S_i<\tau/4$ distinguishes these cases by testing all
estimates against $1-\tau/2$.

Integer normalization ensures $L_i\le S_i^5$ for all large $t$.
Set $A_5=A_5(5)$ before choosing $d$. From~\eqref{eq:three-copy},
\[
 S_i\le6d\,t\sqrt{\log_2t}+O(t),\qquad
 g(S_i)\le(6d+o(1))t.
\]
Choose $6A_5d<1/4$. Polynomially many deterministic queries and the
polynomial-size guesses then give a simulation below $2^{t/2}$ time,
contradicting the hierarchy. Thus infinitely many lengths require
more than $dN\sqrt{\log_2N}$ wires. This is the two-stage
Bathie--Williams proof instantiated at the present scale.

\subsection{Every sufficiently large length}
\label{subsec:ae}

Assume \HTC. Let $n$ denote the input length of the refuter, and
choose
\[
 T(n)=2^{\lceil n^{1/K}\rceil}
\]
for a fixed sufficiently large integer $K$. It is time constructible
and satisfies the refuter's range for all large $n$; finitely many
smaller values can be adjusted. At the sampled lengths $n=t^K$,
write $q(t)=t+O(\log t)$ for a fixed nondecreasing bound on PCP proof
indices and random bits. The two-stage simulator $A_{\mathrm{PCP}}$
just described now runs on arbitrary $z\in\{0,1\}^n$, using the PCP
for $A^T_{\mathrm{FS}}(z)$, with initial PCP soundness $n^{-10}$.
It guesses the two normalized threshold circuits $F,T$, each capped
by $B_d(q)$. The number of guessed bits is polynomial in $t$, with
fixed degree: the two circuit descriptions have that property even
though generating all PCPP constraints can take a larger polynomial
time. Choose $K$ greater than this degree. There are at most $n/10$
guessed bits for all large $n$.

The same error analysis shows
$A_{\mathrm{PCP}}(z)\le A^T_{\mathrm{FS}}(z)$. Its time is
\[
 \poly(n)+\poly(t)2^{(6A_5d+o(1))t}=o(2^t)
 \quad\text{if }6A_5d<1.
\]
The promise holds on every sufficiently large input length, with
$t=\lceil n^{1/K}\rceil$. The refuter therefore produces, for each
large sampled $n=t^K$, a $z_t$ satisfying
\begin{equation}\label{eq:refuter-gap}
 A^T_{\mathrm{FS}}(z_t)=1,
 \qquad A_{\mathrm{PCP}}(z_t)=0.
\end{equation}

It remains to encode both proof-bit functions at the \emph{same}
length. Let $h_0(t)=O(\log t)$ be a fixed canonical header length
and $\ell(t)=O(\log t)$ the uniformly padded proof-bit index length.
Use nondecreasing versions and define
\begin{equation}\label{eq:ae-padding}
 P(t)=h_0(t)+1+q(t)+\ell(t),\qquad
 t(N)=\max\{t:P(t)\le N\}.
\end{equation}
One may take $q(t)=t+C\lceil\log_2(t+2)\rceil$ and analogous
ceiling-log bounds for the other lengths. Thus $P(t+1)-P(t)=O(\log t)$
and $N=q(t)+O(\log t)$ for $t=t(N)$. Use fields
$(\text{header},b,x,j,\text{padding})$, with $|x|=q(t)$ and
$|j|=\ell(t)$. Invalid headers give zero. The $b=0$ branch outputs
the canonical fully accepted PCP proof bit $W_1(t,x)$ for $z_t$ and
ignores $j$; the $b=1$ branch uses the canonical budgeted $F_t$ and
outputs $W_2(t,x,j)$ as above. Both branches ignore padding.
The language is in $\ENP$: the refuter takes $2^{O(t)}$ time with
SAT queries, as do the explicitly expanded canonical searches, and
$t\le N$.

Suppose a circuit of at most $dN\sqrt{\log_2N}$ wires computes this
language at one sufficiently large $N$. Fix the correct header,
padding, and branch bit. For $W_1$ also fix $j$. The resulting
circuits have at most
\[
 dN\sqrt{\log_2N}\le B_d(q(t))
\]
wires for all large $N$, by the factor-two slack in~\eqref{eq:cap}.
The $W_1$ circuit makes a fully accepted $F_t$ exist; this ensures
that the second is a universal PCPP proof circuit for that canonical
$F_t$. Together they make $A_{\mathrm{PCP}}(z_t)=1$, contradicting
\eqref{eq:refuter-gap}. This works separately at every sufficiently
large $N$, proving the worst-case assertion of
Theorem~\ref{thm:lower}. Any sufficiently small positive rational
$d$ with $6A_5(5)d<1$ is admissible.

\subsection{Constant-advantage average-case hardness}
\label{subsec:average}

Assume \HTC. We give the polynomial-parameter version of the proof
in \cite[Section 6]{CLW20}, including the per-term size needed here.
Let the two-query PCPP in \cite[Lemma 3.11]{CLW20} have constants
$0<s_{\mathrm p}<c_{\mathrm p}<1$ and gap
$\Delta=c_{\mathrm p}-s_{\mathrm p}$. Its code bits are parities.
Its constraint count is a power of two, and its proof constructor is
polynomial time. Fix a rational
\begin{equation}\label{eq:delta0}
 0<\delta_0\le\Delta^2/10^6.
\end{equation}
Choose integer exponents $K_1,K_2,K_3$ as specified below, and then
choose the hierarchy's $K$ larger than the fixed degree of all guesses.

The cheating machine $A_{\mathrm{PCPP}}$ obtains the succinct PCP
for $A^T_{\mathrm{FS}}(z)$ and guesses an \emph{ordinary fan-in-two}
circuit $C_1$ of size $t^{K_1}$ as its oracle. The composed verifier
is an ordinary polynomial-size circuit. Apply the two-query PCPP,
with variables $Y\cup Z$, $m=\poly(t)$ constraints, and code word
$Y=\Enc(x)$. Enlarge fixed polynomial bounds on proof lengths and
constraint count uniformly across all candidate $C_1$; the native
power-of-two constraint convention makes binary indexing uniform.
To reach a common larger power of two, repeat every original
constraint equally many times, preserving all satisfaction averages.
The machine guesses, for each proof variable, a rational linear sum
of at most $M=t^{K_2}$ threshold circuits, each with at most $s_1$
wires, and total coefficient absolute value at most $t^{K_2}$.
Coefficients have a fixed $O(\log t)$-bit bound large enough for the
constructions below. These guesses are not assumed to take values
in $[0,1]$.

For a constraint $i$, let $T_{i1},T_{i2}$ be its two (unsigned)
variable functions. Let $F_i$ be its OR polynomial after replacing
a positive literal by $T$ and a negative one by $1-T$. Define
\begin{equation}\label{eq:validity-polys}
 P_{ij}=\begin{cases}
 T_{ij}^2(1-T_{ij})^2,&T_{ij}\text{ corresponds to a }Z\text{ variable},\\
 (\Enc_s-T_{ij})^2,&T_{ij}\text{ corresponds to }Y_s,
 \end{cases}
 \qquad Q_{ij}=T_{ij}^2.
\end{equation}
Estimate expectations of all these polynomials by expanding products
of at most four sum terms. Each expectation is a linear combination
of at most polynomially many expectations of ANDs of at most four
threshold circuits, together with an $O(q)$-wire parity when required.
Every CAPP query has
\begin{equation}\label{eq:four-copy}
 w\le4s_1+O(q).
\end{equation}
The parity is computed as part of the circuit; no enumeration of its
input support is needed. Taking per-product error $t^{-K_3}$ gives
total expectation error $O(t^{4K_2-K_3})=o(1)$ when $K_3>4K_2$.
Increase $K_3$ so all errors below are eventually at most $\delta_0/10$.

Reject unless the estimated average of $P_{ij}$ is at most
$2\delta_0$ and every estimated $Q_{ij}$ is at most $1+\delta_0$.
After this validity test, accept precisely when
\begin{equation}\label{eq:avg-accept}
 \frac1m\sum_i\widetilde{\E}_x F_i
 \ge c_{\mathrm p}-\Delta/2.
\end{equation}
These are the tests of \cite[Section 6.2.2]{CLW20}; their requisite
estimates involve only~\eqref{eq:four-copy}.

Here are the validity estimates explicitly. If the guessed functions
are $[0,1]$-valued and their average $\ell_1$ distance, over occurrences,
from a Boolean PCPP proof is at most $\delta_0$, then the true average
of $P$ is at most $\delta_0$ and every $Q$ is at most one. Thus they
pass. Conversely, passing implies a true average $P\le3\delta_0$
and $\E Q_{ij}\le1+2\delta_0$. Round every $Z$ function to its
nearest Boolean value, and set $Y=\Enc(x)$. For real $u$,
\[
 \min\{u^2,(1-u)^2\}\le4u^2(1-u)^2.
\]
Hence, by Cauchy--Schwarz, the average $L_2$ distance of the guessed
occurrences from the rounded proof is at most
$\sqrt{12\delta_0}$. For each OR polynomial, telescoping its two
arguments, using the Boolean rounded argument and
$\|T_{ij}\|_2\le\sqrt{1+2\delta_0}$, gives
\[
 \|F_i-\widehat F_i\|_1
 \le7\,\E_{j\in\{1,2\}}\|T_{ij}-\widehat T_{ij}\|_2.
\]
For a negated literal use
$\|1-T_{ij}\|_2\le1+\sqrt{1+2\delta_0}$ in the same telescoping;
this gives the constant $7$. If
$A^T_{\mathrm{FS}}(z)=0$, the initial PCP accepts on at most $n^{-10}$
of the $x$'s, so every Boolean PCPP proof has average satisfaction
at most $s_{\mathrm p}+o(1)$. The guessed proof's average is at most
\[
 s_{\mathrm p}+7\sqrt{12\delta_0}+o(1)
 <c_{\mathrm p}-\Delta/2.
\]
Including estimation error preserves the strict inequality. Therefore
$A_{\mathrm{PCPP}}(z)\le A^T_{\mathrm{FS}}(z)$ for all large inputs.

Now fix $\gamma\in(0,1/2)$, and choose the least $k\ge1$ with
\begin{equation}\label{eq:xor-k}
 \eps_k=(1-\delta_0)^{k-1}(1/2-\delta_0)\le\gamma/2.
\end{equation}
Let $c_{\mathrm{STV}}$ be a fixed integer upper bound on the input-length blow-up in the
ordinary-circuit amplification \cite[Lemma 3.9]{CLW20}, originating
in~\cite{STV01}, and put
\begin{equation}\label{eq:kprime}
 k'=\max\{k,c_{\mathrm{STV}}\}.
\end{equation}
Take $s_1=B_{dk}(q)$. Fix
$a_*\ge\max\{5,K_3\}$ large enough for all query descriptions and
precisions, and write $A_5=A_5(a_*)$. Choose a small positive rational
$d$ such that
\begin{equation}\label{eq:avg-time-condition}
 8A_5dk<1.
\end{equation}
Indeed,~\eqref{eq:four-copy} has $4s_1\le8dkq\sqrt{\log_2q}$, so
the machine's time is
$\poly(t)2^{(8A_5dk+o(1))t}=o(2^t)$. The fixed exponent $a_*$ is
chosen to cover the required inverse-polynomial error; the constant
$A_5$ need not equal $A_5(5)$ from the worst-case argument.
All guesses have polynomial length in $t$, so a sufficiently large
$K$ gives the refuter promise. Obtain $z_t$ with
$A^T_{\mathrm{FS}}(z_t)=1$ and $A_{\mathrm{PCPP}}(z_t)=0$.

There are two cases, decidable in $2^{O(t)}$ time with SAT queries by
expanding all PCP random strings.
If no $C_1$ of size $t^{K_1}$ gives a fully accepted oracle, the
canonical fully accepted proof's truth table has ordinary circuit
complexity greater than $t^{K_1}$. The amplification in
\cite[Lemma 3.9]{CLW20} constructs, in $2^{O(q)}$ time from that
truth table, a function $g$ on at most $c_{\mathrm{STV}}q$ inputs
that is $(1/2+t^{-1})$-hard for ordinary size $t^3$, provided
$K_1$ exceeds three times the fixed amplification exponent. The
Boolean simulation of threshold circuits above makes $g$ hard for
the required near-linear threshold wire budget as well.

In the other case take the lexicographically least successful ordinary $C_1$
and a canonical Boolean PCPP proof $\widehat T$, using the proof
constructor separately for each $x$. Set
\[
 f(i,j,x)=\widehat T_{ij}(x),\qquad
 n'=q+\log_2m+1=t+O(\log t).
\]
This function is $\delta_0$-far in $\ell_1$ from every $[0,1]$-valued
linear sum with at most $M/(2m)$ terms, each of at most $s_1$ wires,
and coefficient absolute sum at most $t^{K_2}/(2m)$. To see this,
suppose $D$ is a closer sum and define the guessed function for each
proof variable $X$ by averaging $D(i,j,x)$ over occurrences of $X$.
Empty occurrence lists can be assigned zero. Each resulting function
has at most $M$ terms, no larger per-term wire bound, and no larger
coefficient absolute sum than the allowed cap. For $v_i\in[0,1]$
and $b\in\{0,1\}$,
\[
 \left|\frac1h\sum_i v_i-b\right|=
 \frac1h\sum_i|v_i-b|.
\]
Apply this equality for each $x$ and occurrence list: the average
$\ell_1$ distance to $\widehat T$ remains at most $\delta_0$.
The validity test passes, and the OR polynomials, Lipschitz on
$[0,1]^2$, lose at most $2\delta_0$ average satisfaction relative
to the canonical proof. They therefore satisfy~\eqref{eq:avg-accept},
contradicting $A_{\mathrm{PCPP}}(z_t)=0$.

Apply the \emph{proof of Lemma 3.8 in Appendix A} of~\cite{CLW20},
rather than only its complexity statement. For completeness, its
reconstruction has the following form. Given a circuit approximating
$f^{\oplus k}$, either fix one block and recurse on $k-1$, possibly
negating the circuit, or sample
$\ell=O(n'/(\delta_0\eps_k)^2)$ fixed values of the other blocks.
A Chernoff bound and a union bound over the $2^{n'}$ remaining inputs
make the empirical conditional agreement uniformly accurate. An
affine rescaling of that empirical mean produces a $[0,1]$-valued
linear sum approximating $f$ within $\delta_0$ in $\ell_1$.
Every term is a projection, possibly negated, of the original
approximator. In particular, each term retains its original wire
bound; there are $O_\gamma(t)$ terms, and the sum of coefficient
absolute values is $O_\gamma(1)$, or the weaker sufficient
$O_\gamma(t)$. All coefficients are rational with $O_\gamma(\log t)$
bits because $\delta_0$ and $\eps_k$ are rational. Increasing the
fixed $K_2$ makes these sums fit the preceding caps. Thus
$f^{\oplus k}$ cannot have agreement $1/2+\eps_k$ with the
threshold approximators under consideration.

Finally use a nondecreasing payload bound
\[
 P_\gamma(t)=h_0(t)+
 \max\{k(q(t)+\log_2m(t)+1),\ c_{\mathrm{STV}}q(t)\}
 =(k'+o(1))t,
\]
where polynomial padding of the PCPP supplies a common dyadic
$m(t)$ bound. For each input length $N$ choose the largest $t$ with
$P_\gamma(t)\le N$, and read the active function input from the leading payload bits, ignoring
all other bits. The length $N$ determines $t$; no header is validated and
every word of length $N$ is used. In Case 1, output $g$; in Case 2, output $f^{\oplus k}$. Canonical
searches, case selection, and truth-table amplification all run in
$2^{O(t)}$ time with an NP oracle, so this defines $L_\gamma\in\ENP$.
Set $c_\gamma=dk/k'$. A length-$N$ approximator of at most
$c_\gamma N\sqrt{\log_2N}$ wires has
\[
 s_0\le(dk+o(1))t\sqrt{\log_2t}<B_{dk}(q)=s_1
\]
wires for all large $t$. Averaging over padding gives a fixing with
the same agreement. Both cases yield a contradiction to agreement
greater than $1/2+\gamma$. One can take
\[
 c_\gamma=\Theta\!\left(\frac{1}{A_5k'}\right)
 =\Theta\!\left(\frac{\delta_0}{A_5\log(1/\gamma)}\right)
 \quad\text{as }\gamma\downarrow0.
\]
This proves the average-case assertion of
Theorem~\ref{thm:lower}. The proof uses inverse-polynomial CAPP
precision throughout.

\part{Sparse connectivity}
\section{Sparse determinants with linearly many Boolean wires}
\label{sec:determinants}

The determinant will be a nonnegative connectivity weight after its
permutation terms have been summed. The required tensor representation
must therefore preserve permutation signs as well as matrix entries.
We give the sign construction explicitly; the inverse-pair accounting
is related to the determinant method in~\cite[Section 4]{BCKN15}.

\subsection{A sparse permutation-sign gadget}

Let $E\subseteq[N]\times[N]$ be the potential nonzero positions of
an $N\times N$ matrix, let $M=|E|$, and set
\[
 T_{\mathrm{det}}=N+M+2.
\]
Position bits $z_{ij}$ select one entry per row and column. A legal
selection is a permutation $\sigma$ supported on $E$. We construct
its sign with $O(T_{\mathrm{det}})$ edges and maximum degree
$O(T_{\mathrm{det}}^{2/3})$. The following construction is a
structural identity used under our standing assumption \HTC.

Put $B=\max\{1,\lfloor T_{\mathrm{det}}^{2/3}\rfloor\}$. Rows
having more than $B$ potential positions are heavy; each is a singleton
block. Split every consecutive run of light rows greedily, preserving
row order, so that the total number of potential positions in a block
is at most $B$. Columns are treated independently in their own order.
A zero potential row or column makes the determinant identically zero
and can be handled directly.

If there are $a$ row blocks and $b'$ column blocks, then
\begin{equation}\label{eq:block-count}
 a,b'=O(M/B+1),\qquad ab'=O(T_{\mathrm{det}}^{2/3}).
\end{equation}
Indeed there are at most $M/B$ heavy rows. In one light run, any two
consecutive greedy blocks have total load greater than $B$, so its
block count is at most twice the load divided by $B$, plus one.
There are at most one more light runs than heavy rows. In particular
$a\le4M/B+1$, and the column bound is identical. In a cell
$E_{uv}=E\cap(R_u\times C_v)$, a light side bounds its size by $B$;
if both sides are heavy singletons it contains at most one position.

For a legal permutation write $r(i)$ for its row-block index,
$c(j)$ for its column-block index, and
$n_{uv}=|\{i\in R_u:\sigma(i)\in C_v\}|$. Define
\begin{align*}
 I_R&=|\{i<i':r(i)=r(i'),\ c(\sigma(i))>c(\sigma(i'))\}|,\\
 I_C&=|\{j<j':c(j)=c(j'),\ r(\sigma^{-1}(j))>r(\sigma^{-1}(j'))\}|,\\
 I_{uv}&=|\{i<i':i,i'\in R_u,\ \sigma(i),\sigma(i')\in C_v,
                    \ \sigma(i)>\sigma(i')\}|,\\
 I_\times&=\sum_{u<u',\ v>v'}n_{uv}n_{u'v'}.
\end{align*}
The exact integer identity is
\begin{equation}\label{eq:sign-split}
 \inv(\sigma)=I_R+I_C+\sum_{u,v}I_{uv}+I_\times.
\end{equation}
To prove it, classify a pair $i<i'$ by whether its two row blocks
are equal and whether its two image column blocks are equal. If both
are equal, it is counted in $I_{uv}$. If only the row blocks are
equal, order preservation makes a column inversion equivalent to a
column-block inversion, counted in $I_R$. If only the column blocks
are equal, let $j<j'$ be the two image columns; an inversion is
exactly a reversal of the corresponding row-block order, counted in
$I_C$. If neither is equal, inversions occur precisely in cells with
$u<u'$ and $v>v'$, giving $I_\times$. The four classes are disjoint.

Use constant-state exact-one chains to check row and column selection.
Each light row block reads its at-most-$B$ position bits and contributes
$(-1)^{I_R}$ for that block; a heavy singleton contributes one.
Column blocks similarly contribute $(-1)^{I_C}$. Each cell reads its
position bits, contributes $(-1)^{I_{uv}}$, and outputs one parity bit
\[
 p_{uv}=\sum_{(i,j)\in E_{uv}}z_{ij}\pmod2.
\]
It can also reject repeated selected rows or columns in the cell.
Empty cells output a forced zero bit. A single factor reads the
$ab'$ parity bits and contributes
\begin{equation}\label{eq:cross-root}
 (-1)^{\sum_{u<u',\ v>v'}p_{uv}p_{u'v'}}.
\end{equation}
This gives $(-1)^{I_\times}$. Its degree is bounded by
\eqref{eq:block-count}. Every selected-position bit occurs in only a
constant number of checks. Total port count is
$O(N+M+ab')=O(T_{\mathrm{det}})$; maximum degree is
$O(B+ab')=O(T_{\mathrm{det}}^{2/3})$. Each entry is computed by
polynomially many arithmetic and pair-comparison operations on its
ports. Exact-one states and cell parities are unique on each legal
selection. Equality trees and subdivision preserve these bounds.

For each position multiply the factor
\[
 z_{ij}=0:\ 1,\qquad z_{ij}=1:\ A_{ij}.
\]
Summing over the position bits and the unique states gives precisely
\begin{equation}\label{eq:sparse-det}
 \sum_{\sigma\text{ supported on }E}
    \sgn(\sigma)\prod_iA_{i,\sigma(i)}=\det A.
\end{equation}
The $A_{ij}$ may depend on external Boolean variables. If those
variables and the total scopes of their conditions and numerical
entry factors are $O(T_{\mathrm{det}})$, the combined representation
still has linear edges. Its maximum degree is the larger of the
sign-gadget degree and the external degree. This is a representation
of a determinant inside a larger sum over external choices.

\subsection{Marker Laplacians and positivity}
\label{subsec:matrix-tree}

Fix a simple undirected graph $G=(V,E_G)$, $U\subseteq V$,
$F\subseteq E_G[U]$, and markers $M_0\subseteq U$. Define
\begin{equation}\label{eq:marker-matrix}
 A_{vv}=d_F(v)+[v\in M_0]+[v\notin U],\qquad
 A_{vw}=-[vw\in F]\quad(v\ne w).
\end{equation}
The unselected vertices form an identity block. On $U$, the matrix
is the Laplacian of the graph $H^+$ obtained from $H=(U,F)$ by
adjoining one root $\rho$ and the edges $\rho v$ for $v\in M_0$,
with the root row and column deleted.

For completeness, orient the edges of $H^+$ arbitrarily, and let
$B_\rho$ be its incidence matrix with the root row removed.
Then $A|_U=B_\rho B_\rho^T$, and Cauchy--Binet gives
\begin{equation}\label{eq:matrix-tree}
 \det A=\sum_{\substack{Q\subseteq E(H^+)\\|Q|=|U|}}
            \det((B_\rho)_Q)^2=\tau(H^+).
\end{equation}
A cyclic edge set has dependent incidence columns. An acyclic set
of $|U|$ edges on $|U|+1$ vertices is a spanning tree. Its incidence
minor has determinant $\pm1$, by repeatedly expanding at a nonroot
leaf. Thus every tree contributes one and every other set contributes
zero. Empty $U$ gives determinant one. In particular,
\begin{equation}\label{eq:marker-positive}
 \det A\ge0,\qquad \det A>0
 \ \Longleftrightarrow\
 \text{every component of }(U,F)\text{ meets }M_0.
\end{equation}
If $M_0=\{r\}$ with $r\in U$, positivity is connectivity of $(U,F)$,
and $\det A=\tau(U,F)$ because the root's single incident edge is a
bridge.

The potential positions in~\eqref{eq:marker-matrix} number at most
$n+2m$. For local evaluation, compute every $d_F(v)$ by an
$O(\deg_G(v)+1)$-gate Boolean population counter. One construction
repeatedly uses constant-size full adders to reduce triples in each
binary column: each full adder reduces the total number of live bits
by one, so $O(d)$ suffice for $d$ input bits, followed by addition
of two $O(\log(d+2))$-bit words. Its internal states are functional.
All degree circuits together have $O(n+m)$ wires and gates.
Their output-bit count is $O(n+m)$ as well. A diagonal numerical
factor reads only its degree output, marker bit, selected-vertex bit,
and position bit; an off-diagonal factor reads a position bit and an
edge-selection bit. Total scope length stays $O(n+m)$, including at
vertices of high graph degree.

\subsection{Exact statistics with long binary coefficients}
\label{subsec:weights}

The applications also use a fixed number of statistics
$A_j(y)=\sum_i a_{ji}y_i$ on $p=O(S)$ object bits, with nonnegative
coefficients of at most $H=\lceil S^{9/10}\rceil$ bits. Any fixed
polynomial-time relation on these statistics, including a threshold,
can be added using $O(S)$ wires, degree $O(S^{9/10})$, and a unique
auxiliary extension.

Indeed, group the object bits in blocks of
$B=\max\{1,\lfloor S^{9/10}\rfloor\}$ bits. There are
$h=O(S^{1/10}+1)$ blocks. Every partial or total sum fits in
\[
 q=H+\lceil\log_2(p+1)\rceil+1=O(S^{9/10})
\]
bits. A block factor reads its original bits and checks its output
partial sums. For each statistic, a binary tree of addition factors
combines these outputs. Each addition factor reads three $q$-bit
words, rejects overflow, and enforces exact addition; on genuine
inputs there is no overflow. The root relation reads only the fixed
number of total sums. The total number of state wires is $O(hq)=O(S)$,
whereas the maximum degree is $O(S^{9/10})$. The use of an addition
tree, rather than one factor reading every block sum, is essential to
this degree bound. Local entry arithmetic is polynomial in $S+L$.

For coefficients bounded in value by $S^{O(1)}$, a smaller-block
variant also works: take blocks of $S^{2/3}$ inputs, transmit
$O(\log S)$-bit partial sums, and let a terminal factor read all of
them. Its degree is $O(S^{1/3}\log S)$, and the block degree is
$O(S^{2/3})$. Both encodings are compatible with the determinant
representation. The larger-degree one uses the constant rescaling
\eqref{eq:rescale} before invoking \HTC.

\section{Undirected connectivity: exact sums, optimization, and witnesses}
\label{sec:undirected-proof}

We prove Theorem~\ref{thm:undirected}, always assuming \HTC.
Let $S_G=n+m+2$. The determinant scale in
Section~\ref{sec:determinants} is
\[
 T_{\mathrm{det}}=n+(n+2m)+2=2n+2m+2=O(S_G).
\]
These are distinct parameters; the precision below is measured in
$S_G$.

\subsection{Local candidate systems}

Write $u_v$ for selected-vertex bits and $x_e$ for selected-edge
bits. Requiring selected-edge endpoints to be selected is a constant
scope condition. Incidence degree conditions, including exact degrees
zero, one, and two, have $O(\deg_G(v)+1)$-wire functional chains:
pass the saturated count $\min\{3,\text{count so far}\}$ in two
bits. Domination is checked by an OR chain on each closed neighborhood,
with a total of $O(n+m)$ gates. Vertex cover has the binary constraints
$u_v\vee u_w$ for every edge $vw$. In induced systems introduce the
unique bit $x_{vw}=u_vu_w$ by an AND relation. All such systems have
$O(S_G)$ original and auxiliary bits and total scopes $O(S_G)$.
Global cardinalities and cost relations are the fixed number of
statistics from Section~\ref{subsec:weights}.

For each candidate choose the matrix~\eqref{eq:marker-matrix} and
combine its sparse determinant gadget with these constraints. The
contraction is exactly
\begin{equation}\label{eq:object-sum}
 Z=\sum_{y\text{ satisfying the local and target conditions}}\det A(y).
\end{equation}
The sum over permutation bits is taken inside the sum over $y$; all
other auxiliary states have unique extensions. Table~\ref{tab:objects}
lists the candidates and the determinant value on the objects that
survive the filter. Each single-root row explicitly requires $r\in U$.

\begin{table}[ht]
\centering\small
\renewcommand{\arraystretch}{1.16}
\begin{tabular}{@{}P{.23\linewidth}P{.44\linewidth}P{.25\linewidth}@{}}
\toprule
Problem & Candidate constraints and markers & Surviving weight\\
\midrule
HC / Hamiltonian TSP & $U=V$, all selected degrees two; $M_0=\{r\}$ & $n$ for a Hamiltonian edge set\\
Hamiltonian path & $U=V$, degrees one at distinct $s,t$, two elsewhere; $M_0=\{s\}$ & $1$ for a spanning $s$--$t$ path\\
Simple path & $s,t\in U$, degree one at $s,t$, two at other selected vertices, zero outside; $M_0=\{s\}$ & $1$ for an $s$--$t$ path on $U$\\
Steiner tree & Terminals in $U$, $F\subseteq E_G[U]$, $|F|=|U|-1$; $M_0=\{r\}$ for a fixed terminal & $1$ for a tree $(U,F)$\\
Connected vertex cover & Cover constraints, $F=E_G[U]$, $M_0=\{r\}$ & $\tau(G[U])>0$\\
Connected domination & Domination constraints, $F=E_G[U]$, $M_0=\{r\}$ & $\tau(G[U])>0$\\
Feedback vertex set & $F=E_G[U]$, $M_0\subseteq U$, $|F|+|M_0|=|U|$ & $1$ for each forest with one marker per component\\
\bottomrule
\end{tabular}
\caption{In every row the determinant is zero on candidates failing
the required connectivity condition. Costs and a specified path length
can be added as target relations.}
\label{tab:objects}
\end{table}

Here is the verification of each row. A simple graph of degree two at
every vertex is a disjoint union of cycles. A single marker allows
only the component containing that marker. A surviving cycle on $n$
vertices has exactly $n$ spanning trees, obtained by deleting one
edge. Consequently
\begin{equation}\label{eq:hc-readback}
 Z=n\,\#\mathrm{HC}.
\end{equation}
No orientation or starting-vertex multiplicity is present: the objects
are undirected edge sets. The matrix remains unweighted even when
cost constraints are imposed on these edge sets.

For path candidates, the two odd-degree endpoints must lie in the
same path component, and every other component is a cycle. The
single marker eliminates those cycles; the remaining path has tree
number one. For a specified length $\ell$, impose $|U|=\ell+1$.
For Steiner candidates, connectivity and $|F|=|U|-1$ make the chosen
subgraph a tree. The counted objects are all tree subgraphs containing
the terminal set, including ones with nonterminal leaves. For
connected covers and dominating sets, the determinant is their induced
subgraph's positive tree count, so positivity of~\eqref{eq:object-sum}
is exactly existence.

For FVS, append a root adjacent to every marker. When the determinant
is positive, this augmented graph is connected, and the imposed equation
makes its number of edges equal to its vertex count minus one. It is
therefore a tree and the determinant equals one. Removing its root
leaves the induced forest $G[U]$, with exactly one marker per component.
Conversely every induced forest has such marker choices, all of weight
one. Thus $Z>0$ decides whether a forest of the specified deletion
cost exists. Different marker choices give a positive multiplicity
for the same deletion set.

\subsection{Precision and exact readback}

All matrices in~\eqref{eq:object-sum} have nonnegative determinant by
\eqref{eq:matrix-tree}. Each augmented graph has at most $m+n$ edges,
so its tree count is at most $2^{m+n}$. The numbers of candidates
$(U,F)$, $(U,M_0)$, or $U$ are at most $2^{n+m}$, $3^n$, and $2^n$,
respectively. Uniformly,
\begin{equation}\label{eq:graph-readback-bound}
 0\le Z\le2^{\max\{2n+2m,\,3n+m\}}<2^{4S_G}.
\end{equation}
Take
\begin{equation}\label{eq:graph-precision}
 P_0=4S_G+1,\qquad \ell=1.
\end{equation}
The local signed entries, including $-1$, are legitimate residues in
$\Z/2^{P_0}\Z$. Their intermediate products may cancel; the final
sum is the bounded nonnegative integer in
\eqref{eq:graph-readback-bound}. Its least nonnegative residue is
therefore exactly $Z$. For HC divide that integer by $n$, as in
\eqref{eq:hc-readback}. Even when $n$ is even, no inverse modulo a
power of two is required.

The local candidate constraints, degree counters, and numerical
matrix factors have linear total scope length. The sparse sign gadget
has degree $O(S_G^{2/3})$; adding long cost constraints raises the
degree to $O(S_G^{9/10})$ and keeps $O(S_G)$ wires. Choose
$s=\lceil K S_G\rceil$ with a fixed sufficiently large $K$, and
subdivide edges for bipartiteness. Precision~\eqref{eq:graph-precision},
vertex count, description length, and polynomial-time local evaluation
then satisfy every part of \HTC. One query takes
\begin{equation}\label{eq:one-graph-query}
 \poly(S_G+L)2^{O(g(S_G))}
\end{equation}
time and polynomial space.

\subsection{Optimization, witnesses, and boundary conventions}

Nonnegative edge and vertex costs are attached to the \emph{object}
bits, not to the spanning trees counted by the determinant. For FVS
the objective is the cost of $V\setminus U$, implemented with
complement bits. Use a cost-at-most predicate for minimization and
a cost-at-least predicate for maximization. The sum of the input
costs has $O(n^{9/10}+\log n)$ bits on a sparse graph, and binary
search needs only polynomially many queries. For longest path one
may maximize the number of selected vertices, or an explicitly given
nonnegative selected-edge/vertex weight. The long-coefficient gadget
handles these exact target relations without enumerating cost values.

If a target has $Z>0$, fix its object bits successively. For a bit $y_i$
we have $Z=Z_{y_i=0}+Z_{y_i=1}$, with both summands nonnegative.
Keep any positive child. Apply this only to object bits (and, for
FVS, marker bits), retaining the determinant's permutation sum.
At the end the surviving object satisfies the target and its graph
property. Direct polynomial-time graph and cost checks verify the
returned witness. At most a linear number of additional queries is
needed.

Enumerating possible roots and endpoint pairs costs only a polynomial
factor. For exact counting of unspecified-endpoint undirected paths,
use $s<t$ to count each nontrivial edge set once. Single-vertex paths
are counted separately. For prescribed $s=t$, only the zero-length
single-vertex path is included. Simple undirected cycles require at
least three vertices. An empty terminal set allows the empty Steiner
tree of cost zero in optimization. For counting all nonempty Steiner
tree subgraphs with no terminals, choose the unique root $r=\min U$
and force smaller vertices outside $U$; add the empty tree only if
that convention is requested. With a nonempty terminal set, fix one
terminal as the root, so that no root multiplicity is introduced.
Single-root trees are included. Empty-cover and empty-graph conventions
are checked directly; all nonempty connected sets are covered by the
root enumeration.

The number of queries, including roots, endpoints, binary search, and
self-reduction, is polynomial. Their workspace is reused. Since
$m\le cn$, $g(S_G)=O_c(g(n))$, proving
Theorem~\ref{thm:undirected}.

\section{Directed graphs and sparse-base routing}
\label{sec:directed-routes}

\subsection{The directed marker identity}

Assume \HTC. Let a loopless digraph have selected vertices $U$,
selected arcs $F$ with endpoints in $U$, and markers $M_0\subseteq U$.
Replace~\eqref{eq:marker-matrix} by the outgoing Laplacian
\begin{equation}\label{eq:directed-matrix}
 A^+_{vv}=d^+_F(v)+[v\in M_0]+[v\notin U],\qquad
 A^+_{vw}=-[v\mathbin{\to}w\in F]\quad(v\ne w).
\end{equation}
Its determinant counts spanning arborescences directed toward the
new root $\rho$, after adding $v\to\rho$ for every marker $v$.
Here every selected vertex has one outgoing tree arc and every orbit
reaches $\rho$.

We verify this identity directly. Remove the identity block outside
$U$. Expand each diagonal entry as a choice of one outgoing selected
arc or an available root arc. Fix all these choices, a function
$f:U\to U\cup\{\rho\}$. Its nonroot directed cycles are disjoint.
A nontrivial permutation cycle in a determinant term must be one of
these cycles: its off-diagonal entries require the corresponding
choices of $f$. Conversely any subset of these cycles may be used,
with every other vertex a diagonal fixed point. A length-$d$ cycle
has permutation sign $(-1)^{d-1}$ and $d$ negative entries, so
contributes $-1$. The total for this $f$ is
\[
 \sum_{Q\subseteq\operatorname{cycles}(f)}(-1)^{|Q|}
 =(1-1)^{|\operatorname{cycles}(f)|}.
\]
It is one exactly when every orbit reaches the root. Thus
$\det A^+\ge0$, and it is positive exactly when every selected vertex
can reach a marker along selected arcs. For the converse existence
claim, choose for each vertex an arc decreasing its shortest distance
to a marker or to the new root. Empty $U$ has determinant one.
The exclusion of loops keeps diagonal choices separate from
permutation cycles.

The potential matrix positions number at most $n+m$, and outgoing
degree counters and entry scopes are $O(n+m)$. Combining this
identity with the sparse sign gadget obeys the same wire, degree,
entry-cost, and ring bounds as before.

\subsection{Directed applications}

\begin{proof}[Proof of Theorem~\ref{thm:directed}]
For directed HC, require indegree and outdegree one at every vertex
and fix one marker $r$. Candidates are disjoint directed cycles.
A spanning cycle has exactly one rooted arborescence in the augmented
graph: replace $r$'s outgoing cycle arc by its root arc, and keep the
other arcs. Any additional cycle misses the marker and has determinant
zero. The final sum is therefore the directed HC count itself; it
is not divided by $n$.

For a directed $s$--$t$ path, $s\ne t$, require selected indegree/outdegree
$(0,1)$ at $s$, $(1,0)$ at $t$, $(1,1)$ at other selected vertices,
and $(0,0)$ outside $U$. Use $t$ as the sole marker. The path contributes
one, and any extra directed cycle makes the determinant zero. This
proves exact counting with a prescribed length and longest-path
optimization. Unspecified endpoints are enumerated as \emph{ordered}
pairs. A single-vertex path is treated separately. A two-vertex directed
cycle with opposite arcs is a valid HC under these conventions.
Empty graphs and the one-vertex loopless case are checked directly.

For directed Steiner tree toward a root $r$, require $r\in U$,
outdegree zero at $r$, outdegree one at each other selected vertex,
and inclusion of all terminals. With marker $r$ the determinant is
one precisely on arborescences toward $r$. A feasible rooted-reachability
subgraph contains a suitable arborescence, obtained by selecting arcs
that decrease distance to the root; extra branches not needed for the
terminals can be deleted. With nonnegative costs this restriction
preserves the optimum. Reverse the input for the usual version with
paths from the root to the terminals. The same network exactly counts
arborescent subgraphs of the specified direction. Even when terminals
consist only of the root, the counting problem includes all such
subgraphs; only the optimum can immediately be the single root.

For a strongly connected selected subgraph, fix $r\in U$ and multiply
two determinant gadgets on the \emph{same} selected-arc bits: one for
$F$ and one for its reverse. The second matrix's diagonal uses the
original indegrees; it is recomputed rather than obtained by merely
transposing the original matrix. The product is positive exactly when
every selected vertex can reach $r$ and can be reached from $r$.
This is strong connectivity. The product stays inside the common
sum over candidates. It is not the product of two separately
contracted sums. For a spanning subgraph use $U=V$. For a terminal
version use $T\subseteq U$ and a fixed $r\in T$ when $T$ is nonempty.
For an empty terminal set use the requested empty-subgraph convention,
or enumerate roots to cover nonempty solutions.

Each arborescence is a subset of at most $m+n$ potential augmented
arcs. At most two determinant factors and at most $2^{n+m}$
candidates give the uniform bound
\[
 0\le Z\le2^{3(n+m)}<2^{4(n+m+2)}.
\]
Use $P_0=4(n+m+2)+1$, $\ell=1$, and a fixed multiple of $n+m+2$
for the \HTC\ scale. Costs use the exact statistics gadget; objective
and witness bits are recovered by the same positive-sum self-reduction
as in the undirected case. All query counts are polynomial. This
proves the theorem.
\end{proof}

\subsection{Closed walks, required edges, and metric TSP}

\begin{proof}[Proof of Theorem~\ref{thm:routing}]
A nonempty closed walk determines a connected even-degree edge
multiset. If an edge has multiplicity at least three, delete two of
its copies. Its support remains unchanged, degrees stay even, and
required-edge and terminal constraints remain satisfied. The cost
does not increase. Repeating gives an optimum with multiplicities
\begin{equation}\label{eq:walk-multiplicity}
 \mu_e\in\{0,1,2\}.
\end{equation}
Conversely a connected even-degree multigraph has an Euler tour:
follow unused edges to a closed trail, and splice in further closed
trails wherever unused edges remain. Connectivity and parity ensure
this process exhausts every edge. A selected single vertex without
edges represents its zero-length closed walk.

Encode $\mu_e=a_e+2b_e$ with Boolean bits and forbid $a_e=b_e=1$.
Define its support $f_e=a_e\vee b_e$ uniquely. Required edges have
$f_e=1$. Select $U$ containing the terminals and every support
endpoint; the single-marker determinant on the support graph
$(U,\{e:f_e=1\})$ enforces connectivity. Extra isolated selected
vertices are thereby eliminated except for the valid singleton case.
For nonempty requirements fix
$r\in T\cup V(R)$ as the marker and require $r\in U$. If a depot is
specified, include it in $T$. Empty requirements allow an empty walk
of cost zero and are handled directly.

A parity chain at vertex $v$ requires
\[
 \sum_{e\ni v}a_e\equiv0\pmod2,
\]
which is exactly the parity of its multidegree. The chains and the
multiplicity/support relations together have $O(n+m)$ wires. The
cost is the exact statistic
\[
 \sum_e w_ea_e+2w_eb_e.
\]
Its coefficient width is at most one more than the input width;
this constant change is covered by the fixed-scale long-coefficient
gadget. The determinant counts trees of the \emph{simple support}
graph, as a positive feasibility weight.

There are at most $2^n3^m$ candidates $(U,\mu)$, and at most $2^m$
support spanning trees for a single marker. Thus
\[
 0\le Z\le2^n6^m<2^{4(n+m+2)}.
\]
The same precision and \HTC\ scale apply. Binary search on cost,
followed by fixing $a_e,b_e,u_v$ while preserving positive sums,
recovers an optimal even connected multiset. Its Euler tour has at
most $2m$ edge traversals. With a depot, rotate the tour to start
and end there.

Finally suppose the base graph is connected, $T=V$, and $R$ is empty.
Expand any Hamiltonian tour of its shortest-path metric into base-graph
shortest paths. This gives a closed spanning walk of the same cost.
Conversely, visit vertices in their order of earliest appearance along a
spanning walk and shortcut repeats in the metric. The triangle
inequality gives a Hamiltonian metric tour of no greater cost.
Therefore the optima are equal, and shortcutting an optimal returned
walk gives an optimal metric tour. Zero costs can give a pseudometric;
the same inequalities still hold. Distances are needed only for this
polynomial-time output step; a dense metric closure is not part of
the tensor network. The proof of~\eqref{eq:graph-time} follows.
\end{proof}

\section{Encoding boundaries and finite checks}
\label{sec:checks}

\subsection{Precision and representation boundaries}

At $S=\Theta(n)$, the high-precision error in
Theorem~\ref{thm:capp} is still much larger than $2^{-n}$. Its
absolute-error guarantee does not distinguish zero accepting inputs
from one. Exact modes require their weight-cost budget or a supplied
short factorization. A general polynomial-bit equality constraint on
$n$ input bits need not satisfy those exact budgets.

The sum in~\eqref{eq:budget} is a sufficient condition for the
constructions above; no necessity is claimed for circuits whose
gates share realizations. The seed estimate used by
the algorithm is $R=O_{a,c}(S^{9/10})$, measured in the scale $S$.
Exact-only assignments have $R=0$; hashed assignments use their
pointwise error bounds for the precision comparison.

A ring calculation clarifies the cost of a particular congruence
filter. In $R=(\Z/2^P\Z)[V]/(V^\ell)$ let $\mathfrak m=(2,V)$.
If units $\zeta_j\in R$ and coefficients $c_j\in R$ satisfy
\[
 \sum_jc_j\zeta_j^u=2^e[u=0],\qquad 0\le u<M,\quad e<P,
\]
then $e\ge M-1$. Every unit is one modulo $\mathfrak m$; the
$(M-1)$st forward difference of the left side belongs to
$\mathfrak m^{M-1}$, while that of the right side is
$(-1)^{M-1}2^e$. The constant coefficient of
$\mathfrak m^{M-1}$ is divisible by $2^{M-1}$, forcing the bound.
Filters of this specified unit-power form therefore consume precision
proportional to their modulus range.

\subsection{The fixed-width counting interface}

The local condition ``two selected incident edges at every vertex''
recognizes a spanning 2-factor. Two disjoint triangles satisfy it and
have no HC. Joining them by a bridge and leaving the bridge unselected
gives the same example on a connected input. Rooted consistent cuts
give exactly $2^{\cc(F)-1}$ extensions and filter disconnected
2-factors modulo two. Distinct Hamiltonian solutions in one weight
layer can still cancel. The deterministic proofs use nonnegative
determinant sums instead.

Section~\ref{sec:lean-connectivity} supplies a linear fixed-width encoding
of each Cut\&Count weight layer, giving randomized decision and witnesses
under \HLean. Numbering-based certificates with $\Theta(n\log n)$
Boolean variables would give only $2^{O(n\sqrt{\log n})}$ even under
the manuscript's linear-density bound. The digit grouping avoids that
cost for fixed polynomially bounded coefficients. The deterministic
decision, full integer counting, and long-cost results of
Sections~\ref{sec:determinants}--\ref{sec:directed-routes} still use \HTC;
the new reduction does not establish those statements from \HLean.

One \HBB-only example can be verified directly. On a sparse graph,
exact chromatic number and an optimal coloring are computable,
conditional on \HBB, in
$\poly(n+L)2^{O(n\log\log n/\sqrt{\log n})}$ time and polynomial
space. Take $D=\lceil\log_2(n+2)\rceil^2$ and let $H$ contain
vertices of original degree greater than $D$. Then
$|H|=O(n/\log^2n)$. Compute $\chi(G[H])$ by enumerating all
$|H|^{|H|}$ color labelings. If it is at least $D+1$, an optimal
coloring of $H$ extends greedily to the remaining vertices and is
optimal. Otherwise try $q\le D+1$ and every $q^{|H|}$ coloring of
$H$. Represent each remaining color in $O(\log q)$ bits, reject
codes at least $q$, and test each remaining or crossing edge with an
$O(\log q)$-gate inequality circuit. The functional Tseitin 3-CNF
has $O(n\log\log n)$ variables and clauses. Add one fixed-zero dummy
bit per clause, plus twenty, to make density less than two while
preserving the count. Apply \HBB\ with fixed $k=3$ and density two;
all enumeration factors are absorbed in the stated exponent. Fixing
original color bits recovers an optimal extension. Empty graphs and
the case $q=1$ are checked directly. This is a routine
composition of degree splitting, binary comparison, and \HBB.

\subsection{Budget scales in lower-bound simulations}

All \HTC\ calls in Section~\ref{sec:lower} go through
Theorem~\ref{thm:capp}, whose scale is the actual $S$ of the query
circuit; the proofs use only upper bounds on the resulting running time,
so no lower bound on the size of a guessed circuit is needed.

For average-case amplification, the input length grows by $k$ and the
analysis-query exponent is proportional to $ck$. Condition
\eqref{eq:avg-time-condition} permits constant $k$ at a fixed positive
wire constant. Increasing CAPP precision alone does not remove that
time constraint.

\subsection{Classical sparse baselines}

The sparse reductions for the usual ETH comparison have explicit
linear budgets. Sparsification~\cite{IPZ01} gives a 3-CNF with $v$
variables and $a=O(v)$ clauses. The standard vertex-cover construction
has a variable edge and a clause triangle, with each clause vertex
joined to its literal's variable vertex. Its budget is $v+2a$;
there are $2v+3a$ vertices and at most $v+6a$ edges. The forced
one-per-variable and two-per-triangle choices make an unchosen clause
vertex imply a true literal. The standard dominating-set construction
has triangles $(t_i,f_i,z_i)$ and a vertex for each clause joined to
its literals, with budget $v$. The private $z_i$ forces one selection
per triangle; choosing $z_i$ can be replaced by a literal. It has
$3v+a$ vertices and at most $3v+3a$ edges.

From vertex cover, make a new triangle on each retained edge by adding
one vertex. Deleting a cover leaves a forest; a selected new vertex
in an FVS can be replaced by an endpoint without creating a cycle.
Adding a universal root and $k+2$ private leaves gives connected
vertex cover with budget $k+1$. For a dominating-set instance use
selectors $a_u$, elements $b_v$, edges $a_ub_v$ for $v\in N[u]$,
and a root joined to every selector. Root plus $k+2$ private leaves,
with budget $k+1$, gives connected domination. Selected elements can
be replaced by adjacent selectors because the root already dominates
all selectors. Alternatively declare the root and all elements Steiner
terminals: a tree using $s$ selectors has $n+s$ edges, so budget
$n+k$ is equivalent to $s\le k$. These graphs have linear vertex
and edge counts, and edge-to-vertex ratios at most four. They give
the standard sparse ETH baselines for these problems. The algorithms
in this paper use their explicit \HTC\ encodings.

\subsection{Ancillary computations}

The directory \ancfile{anc/} contains finite checks of the reductions.
Each check prints one summary line and writes its details to
\ancfile{anc/results/}. The counting checks use
integer or rational arithmetic; the floating-point threshold scan is
a diagnostic. Full exhaustive families and exhaustive
checks on sampled supports are distinguished in Table~\ref{tab:checks}.
The runner executes the checks serially.

\begin{table}[ht]
\centering\small
\renewcommand{\arraystretch}{1.12}
\begin{tabular}{@{}P{.37\linewidth}P{.55\linewidth}@{}}
\toprule
Check & Actual finite coverage\\
\midrule
\ancfile{finite_audit.cpp} & All permutations and interval partitions through $N=7$: 21,413,481 sign identities; 254,254 undirected marker cases through $n=5$; 333,544 directed cases through $n=4$\\
\ancfile{p6_sign_check.py} & 16,362 supported permutation identities; 350 sampled-support factor sums, including 218 explicit parity-wire sums\\
\ancfile{p7_p8_filter_check.py} & 120 sampled undirected graphs, 76 directed graphs; object sums and 68 metric-TSP comparisons\\
\ancfile{e2e_hc_network.py} & Nine graph/blocking combinations, including disjoint triangles; modular readback and integer division\\
\ancfile{check_gadgets.py} & 9,000 arithmetic replays of sliced threshold/equality/modular formulas, residue hashes, and small modular-function enumeration; no tensor contraction\\
\ancfile{check_G1_tensor.py} & 3,816 consistent-input factor contractions and 411 independently assigned slice-copy cases, including full auxiliary code domains\\
\ancfile{check_G3_tensor.py} & 136,896 residue-chain factor contractions on sampled instances, all their inputs and 24 primes; invalid residue codes included\\
\ancfile{protocol_exact_check.cpp} & Two parties, eight-bit inputs, 2,048-prime family: 3,346,110 nontrivial state/threshold classes\\
Explicit protocol tensor checks & Local tensors, equality subdivision, generic exact elimination, early stops and wrong-hash seeds; one extension for every original input\\
\bottomrule
\end{tabular}
\caption{Validation of the reductions. Auxiliary engine, block-count,
and unused-input checks and commands appear in \ancfile{anc/README.md}.}
\label{tab:checks}
\end{table}

The program \ancfile{hc_accounting.py} gives approximately $42.8S_G$ to
$52.4S_G$ bipartite wires on its five listed sparse fixtures, with
$P_0=4S_G+1$. These are finite examples; the asymptotic bounds follow
from the port counts proved above. The computations do not execute
the algorithm asserted by \HTC. Source hashes, commands, logs,
and scopes accompany the programs.
The arithmetic replay in \ancfile{check_gadgets.py} does not test
tensor-level extension counts, invalid auxiliary codes, or independently
assigned copies; its modular enumeration does not implement the general
integer-representation construction.
The tensor checks \ancfile{check_G1_tensor.py} (sliced gadgets) and
\ancfile{check_G3_tensor.py} (residue chains) cover the stated full-code
domains. In the sliced-gadget check, independently assigned slice copies test the local gadget's
extension count; the circuit's equality trees impose consistency of
those copies. The ancillary numerical threshold scan is a finite scan,
and the eventual inequalities are justified by the asymptotic proofs.

For Section~\ref{sec:lean-connectivity}, \ancfile{cnf_circuits.py}
constructs the functional 3-CNF gates and grouped weighted sums.
The \ancfile{check_cnf_*.py} programs compare literal-level exact
counts with direct graph and arithmetic references. They test
weight-layer cut multiplicities, odd-branch recovery, polynomial-cost
optimization, marker isolation, terminal-selection units, and the
unsigned modular expansion. Counterexamples include a feasible but
suboptimal TSP output and a missing-terminal directed Steiner formula.
The arithmetic checks include fixed inputs across a radix boundary;
large gate-accounting cases are distinguished from actual CNF solving.
The precise graph ranges, commands, and elapsed times are recorded in
\ancfile{anc/README.md} and \ancfile{anc/results/manifest.json}.
These checks use elementary exact counters, not the fast algorithm
assumed by \HLean, and do not formalize the reductions in Lean.

\noindent\begin{minipage}{\linewidth}
\subsection*{Use of AI}
The proofs and the text were produced with AI systems (GPT-6.1 Sol,
GPT-6 Astra, Claude Opus 5.5, and Claude Fable 5.1) under the author's
direction. They were checked by independent reviews with AI models and by the
computations in the ancillary files. No human expert has reviewed
them. The author takes responsibility for the content.
\end{minipage}\par

\begingroup
\small
\raggedright
\bibliographystyle{plain}
\bibliography{references}
\endgroup
\end{document}
